% GUID v5: Quantum Gravity. A complete, standalone paper; it cites DMP 0.9.9 and the physics literature only.
% Source of the paper. Build: see README.md in this folder (pdflatex twice; no BibTeX, the bibliography is inline).
% Licence: Apache License 2.0.
\documentclass[11pt,a4paper]{article}
\usepackage[T1]{fontenc}
\usepackage[utf8]{inputenc}
\usepackage{lmodern}
\usepackage[margin=2.5cm]{geometry}
\usepackage{amsmath,amssymb,amsthm}
\usepackage{booktabs,tabularx,array,longtable}
\usepackage{enumitem}
\usepackage[hidelinks]{hyperref}
\usepackage{url}
\usepackage{xcolor}

\newtheorem{theorem}{Theorem}
\newtheorem{proposition}{Proposition}
\newtheorem{lemma}{Lemma}
\newtheorem{corollary}{Corollary}
\theoremstyle{definition}
\newtheorem{definition}{Definition}
\newtheorem{hypothesis}{Hypothesis}
\theoremstyle{remark}
\newtheorem{remark}{Remark}
\newcolumntype{L}{>{\raggedright\arraybackslash}X}

\setlength{\emergencystretch}{3em}
\newcommand{\chk}[1]{\textsf{[#1]}}
\newcommand{\tagD}{\textsf{(D)}}
\newcommand{\tagH}{\textsf{(H)}}
\newcommand{\tagP}{\textsf{(P)}}
\newcommand{\ket}[1]{\lvert #1\rangle}
\newcommand{\bra}[1]{\langle #1\rvert}
\newcommand{\Tr}{\operatorname{Tr}}
\newcommand{\dd}{\mathrm{d}}
% Version 5.1: new text is set in this colour and tagged [v5.1].
\definecolor{newc}{RGB}{0,70,150}
\newcommand{\nw}[1]{{\color{newc}#1}}
\newcommand{\vtag}{{\color{newc}\textsf{\scriptsize[v5.1]}}}
\newenvironment{vnew}{\color{newc}}{}
\theoremstyle{definition}
\newtheorem*{hypothesisprime}{Hypothesis 1$'$}
\newtheorem*{hypothesisprime6}{Hypothesis 6$'$}

\title{GUID v5: Quantum Gravity\\[0.6ex]
\large A classical existence frame carrying quantum information: what the axioms of the Dark Matter Protocol fix about gravity,
what they leave to quantum theory, and the experiments that decide}
\author{Steven Hrovat\thanks{Independent. Contact: \href{mailto:hrovat.steven@googlemail.com}{\nolinkurl{hrovat.steven@googlemail.com}}.}
\and Sirius (Claude, Anthropic)\thanks{An AI system by Anthropic, working as the coordinator of the first author's projects. See
``Contributions and use of AI''.}}
\date{Version 5.1, draft of 31 October 2026}

\begin{document}
\maketitle

\begin{abstract}
\noindent The Dark Matter Protocol (DMP~0.9.9) is an information-theoretic axiom system: every act needs a Duality of two independent
witnesses, one owning trust and one owning information; trust is a field; and every act leaves a record that its witnesses read in full
while any other reader sees only counts. GUID reads these axioms as physics; this paper states its account of quantum gravity. \nw{The
trust field is the gravitational plane, a tetrad shaped by the counts of existence and information but separate from both; existence is
the rest mass the Higgs vacuum grants,} information the quantum content and a scalar that sets the fine-structure constant. From the axioms
we derive universal free fall, light bending with $\gamma=1$ and, by the no-broadcasting theorem, that the \nw{plane} is classical once the
start is complete. Supplied with the base configuration currently assumed for our universe, DMP hosts this theory. Consistency with quantum
mechanics requires the \nw{plane} to diffuse, and the axioms' record-and-recompute structure is exactly a measurement-and-feedback channel
of the Kafri-Taylor-Milburn and Tilloy-Di\'osi type. We also treat horizons, black holes, holographic codes, 2+1 dimensions and cosmology,
\nw{from a start that splits the strong from the electroweak force to the electroweak split, where the Higgs vacuum and $\alpha$ appear}.
The theory lets remnants of interactions be trapped in the gravitational field. It predicts no gravitationally mediated entanglement,
spacetime diffusion tied to decoherence, $\gamma=1$ and universal free fall. Made parameter-free by its own principles, it gives the
Di\'osi-Penrose rate: the route is new, the physics belongs to the classical-quantum family. All numbers come from one public script of
\nw{123} checks. We claim that the construction exists and reproduces, not that it is correct. \nw{Text new in version 5.1 is set in
blue.}
\end{abstract}

\noindent\textbf{Keywords.} quantum gravity; classical-quantum dynamics; gravitationally mediated entanglement; no-broadcasting theorem;
tetrad gravity; spacetime diffusion; black-hole thermodynamics; holographic codes; information-theoretic axioms.

\noindent\textbf{Subject.} General Relativity and Quantum Cosmology (gr-qc); related: Quantum Physics (quant-ph), Information Theory
(cs.IT).

\tableofcontents

\section{Introduction}\label{sec:intro}

\subsection{The problem}

General relativity describes gravity as the curvature of a classical spacetime metric $g_{\mu\nu}$ sourced by energy and momentum,
$G_{\mu\nu}+\Lambda g_{\mu\nu}=(8\pi G/c^4)T_{\mu\nu}$. Quantum theory describes every other interaction as a quantum field. The two
meet in three places, and in each the meeting is not understood.

\emph{The ultraviolet.} Quantising the metric perturbatively around flat space, $g_{\mu\nu}=\eta_{\mu\nu}+\kappa h_{\mu\nu}$ with
$\kappa=\sqrt{32\pi G}/c^2$, gives a massless spin-2 particle, the graviton, whose coupling has the dimension of a length. The theory is
not renormalisable: pure gravity diverges at two loops (Goroff and Sagnotti~\cite{goroff}), and with matter already at one loop ('t~Hooft and
Veltman~\cite{thooftveltman}). Below the Planck energy $E_P=\sqrt{\hbar c^5/G}\approx1.22\times10^{19}$\,GeV the theory is nevertheless a
reliable effective field theory (Donoghue~\cite{donoghue}); its leading quantum correction to the Newtonian potential is
\begin{equation}\label{eq:eft}
V(r)=-\frac{Gm_1m_2}{r}\left[1+3\frac{G(m_1+m_2)}{rc^2}+\frac{41}{10\pi}\frac{G\hbar}{r^2c^3}\right],
\end{equation}
where the middle term is classical and the last is quantum~\cite{bdh}. At a distance of one metre the quantum term is about
$3\times10^{-70}$ of the first \chk{Q1}. What completes the theory above $E_P$ is the subject of string theory, loop quantum gravity, asymptotic safety, causal dynamical
triangulations, causal sets and others (Section~\ref{sec:landscape}).

\emph{Black holes.} A black hole radiates thermally at the Hawking temperature $k_BT_H=\hbar c^3/8\pi GM$~\cite{hawking75} and carries the
Bekenstein--Hawking entropy $S=k_Bc^3A/4G\hbar$~\cite{bekenstein73}. Whether the radiation carries the information that fell in, and how,
is the black-hole information problem~\cite{page93,penington,aemm}.

\emph{Whether to quantise at all.} Nothing in experiment yet requires the gravitational field itself to be quantum. Feynman's argument for
quantising it, the Eppley--Hannah thought experiment~\cite{eppley} and the Page--Geilker experiment~\cite{pagegeilker} rule out specific
ways of coupling a classical metric to quantum matter, not all of them. A consistent classical-quantum dynamics exists and is testable
(Oppenheim~\cite{oppenheim}; Oppenheim, Sparaciari, \v{S}oda and Weller-Davies~\cite{osswd}), and the question has become experimental:
if gravity can entangle two masses through local interaction, it is not classical (Bose et al.~\cite{bose}; Marletto and Vedral~\cite{mv};
Galley, Giacomini and Selby~\cite{ggs}).

\subsection{What this paper does}

GUID, the General Unified Information Duality, is the reading of an information-theoretic axiom system, the Dark Matter Protocol
(DMP~0.9.9~\cite{dmp}), as physics. DMP was written to govern trust in computer systems. Its first axiom asks that every act be passed by two
independent witnesses, one owning trust and one owning information; its second calls trust ``gravitational'', a field with sources,
depths and a freshness that decays; its third asks that every act leave a record that its witnesses read in full while every other reader
sees only counts, and that nothing about information be derivable from trust except that it exists. The structural resemblance to physics
is deliberate: a field that sources from mass, boundaries that contain, a floor at which only entropy leaves.

This paper collects what follows when that resemblance is taken literally, and states it as one account of quantum gravity. The account
has three parts. \emph{What the axioms fix} (Section~\ref{sec:fix}): with trust read as existence, gravity is carried by a frame that falls
universally, does not mix with the information half, and is classical, because the axioms make every one of its values a record that two
witnesses hold in full; it therefore mediates no entanglement. Supplied with the base configuration currently assumed for our universe,
DMP hosts this theory: its own shared structure carries the classical frame, and its information half the quantum content. \emph{What consistency then demands} (Sections~\ref{sec:cq} and~\ref{sec:exp}): a classical frame coupled to quantum matter must
decohere that matter and must itself diffuse, with a trade-off between the two that experiments bound from both sides; the axioms' own
architecture, record then recompute, is exactly the measurement-and-feedback channel that realises this. \emph{What it says elsewhere}
(Sections~\ref{sec:thermo} to~\ref{sec:matter}): horizon thermodynamics, black holes, holographic codes, lower dimensions, cosmology and
the matter content, each tested against measurement, including three routes by which the theory traps remnants of interactions in the
gravitational field and what each predicts for dark matter (Section~\ref{sec:darkmatter}).

\paragraph{Result in one paragraph.} In GUID, quantum gravity is not a quantisation of geometry. It is a classical existence frame, the
tetrad of general relativity, over quantum information, and DMP, supplied with the base configuration currently assumed for our universe,
hosts it. Free fall is universal, $\gamma=1$, gravitational waves are classical waves of the frame that travel at $c$, and the frame is
classical once the start, inflation together with the decay of the oscillons it ends in, is complete. The theory is on the classical side of
the Galley--Giacomini--Selby no-go theorem, with Jacobson's thermodynamic gravity, Oppenheim's classical-quantum gravity and the
Kafri--Taylor--Milburn and Tilloy--Di\'osi channel models. It lets remnants of interactions be trapped in the gravitational field, which gives
it two live candidates for dark matter. It is falsifiable by a positive test of gravitationally mediated entanglement, by a test of the LOCC
inequality of Lami, Pedernales and Plenio, and by the decoherence-diffusion trade-off. A theory in which the geometry itself is quantum can
still run on DMP as content, so the axioms do not forbid quantum gravity in nature; such a geometry would not be DMP's frame.
\nw{\vtag{} In version 5.1 gravity is the separate plane that existence and information shape; the Higgs vacuum is the existence that
separated out when the information field split from the electroweak field; and causality begins with the split of the strong force from
the electroweak force.}

\paragraph{What we claim.} We claim that the derivations are as written, that every number reproduces from one public script
(\texttt{docs/quantum-gravity/checks.py}, \nw{123 of 123}) and from the holographic-code assembly (\texttt{docs/qg-verdict/assembly.py},
18 of 18), and that the verdicts follow from the measurements cited. We do not claim that the account is correct, and we name in
Section~\ref{sec:break} where we expect it to break.

\paragraph{What is new and what is not.} Every physical prediction of this paper already belongs to an existing programme: classical
gravity coupled to quantum matter and its decoherence-diffusion trade-off~\cite{oppenheim,osswd}, gravity as a measurement-and-feedback
channel~\cite{ktm,tilloydiosi}, gravitationally mediated entanglement as the test~\cite{bose,mv,ggs}, horizon
thermodynamics~\cite{unruh,hawking75,jacobson95}, Christodoulou's mass formula~\cite{christodoulou}, holographic codes with a classical
centre~\cite{harlow}, and the dark-matter routes of Section~\ref{sec:darkmatter}. Made parameter-free by its own principles, GUID lands on
the Di\'osi--Penrose rate or on a refuted value (Section~\ref{sec:bridges}). GUID is therefore not a new theory of quantum gravity; it is a
member of the classical-quantum family. What is new is the route and four statements: (i) a derivation of the classical side of the open
question from an information principle, namely that a structure every party can audit in full cannot be quantum (Lemma~\ref{lem:classical}),
where the literature postulates it or leaves it to experiment; (ii) the identification of DMP's record-and-recompute architecture with
measurement-and-feedback gravity (Proposition~\ref{pr:ktm}); (iii) one rule of its own, that DMP's balance holds the decoherence-diffusion
trade-off at equality, so that a measured gravitational decoherence would fix the spacetime noise (Section~\ref{sec:bridges}); (iv) the
converse lesson for the design of information systems, that a control structure every party can audit in full is classical while the data
it governs need not be. All four rest on Hypothesis~\ref{h:exist}, which is an identification, not a derivation. We regard the paper as a
contribution to the foundations of physics: an information-theoretic principle that selects classical-quantum gravity, with falsifiers of
its own.

\begin{vnew}
\paragraph{Version 5.1.}\vtag{} This version folds in four readings of the first author, each checked against measurement and fitted into
the section where it belongs. Text that is new in version 5.1 is set in this colour and tagged \vtag. (i) Gravity is a separate plane,
shaped by the existence and the information it contains but never either of them (Hypothesis~1$'$, Section~\ref{sec:fields}); this
removes a tension the earlier hypothesis had with free fall, which weighs all energy. (ii) The Higgs vacuum is the existence that separated
out when the information field, the photon and with it $\alpha$, split from the electroweak field, about $10^{-11}$\,s after the start
(Section~\ref{sec:higgsplit}). (iii) $\alpha$ was never measurably one: it was born near $1/128$ at that split (Section~\ref{sec:alpha}).
(iv) Causality begins with the split of the strong from the electroweak force, within the start; a unification, if there was one, lies
before time (Hypothesis~6$'$, Section~\ref{sec:cosmo}). The ledger, the predictions and the checks (\chk{Q20}) follow.
\end{vnew}

\paragraph{Contributions.}
\begin{enumerate}[nosep]
\item A formal statement of DMP's axioms in the form physics needs, and the dictionary that reads them as two fields (Sections~\ref{sec:dmp},
\ref{sec:fields}).
\item Derivations from the axioms of universal free fall, $\gamma=1$, the classicality of the frame, and the statement that DMP, supplied
with the base configuration of our universe, hosts the proposed theory (Section~\ref{sec:fix}).
\item The identification of the axioms' record-and-recompute structure with the measurement-and-feedback channel of classical-quantum
gravity, the balance rule that holds its trade-off at equality, the two bridges that fix its constant, and the predictions that follow
(Sections~\ref{sec:cq},~\ref{sec:exp}).
\item A consistent treatment of horizons, black holes, holographic codes, 2+1 gravity, cosmology and matter in that account
(Sections~\ref{sec:thermo}--\ref{sec:matter}), with a comparison to the main approaches to quantum gravity (Section~\ref{sec:landscape}).
\item Three routes by which the theory traps remnants of interactions in the gravitational field, with what each predicts for dark matter
(Section~\ref{sec:darkmatter}).
\item A table of predictions with the experiments that test them, and a ledger of every statement with its verdict
(Sections~\ref{sec:predictions},~\ref{sec:ledger}).
\end{enumerate}

\subsection{How to read this paper}

Every step carries one of three marks. \tagD{} follows from the text of DMP~0.9.9 (commit \texttt{4400cce} of its repository~\cite{dmp}).
\tagH{} is a hypothesis of GUID: an identification of a DMP word with a physical quantity, which no derivation supplies and which the
paper states once, in Section~\ref{sec:fields}, and uses by name. \tagP{} is established physics, with its source. Every statement that
meets data receives one of five verdicts: \emph{derived} (it follows from the axioms and the hypotheses named), \emph{holds} (it agrees with
measurement and follows from the steps given), \emph{holds as a reading} (it agrees, but adds nothing a measurement could tell apart from the
established account), \emph{refuted} (a measurement or calculation contradicts it), \emph{untested} or \emph{live} (no measurement yet reaches
it, or one is under way). A label such as \chk{Q5} names a check in \texttt{docs/quantum-gravity/checks.py}; \chk{A3} names check~3 of the
assembly script. Units are SI unless stated; where we set $G=c=1$ or $\hbar=c=k_B=1$ we say so. The metric signature is $(-,+,+,+)$.

\section{The axiom system}\label{sec:dmp}

DMP opens with undefined words, then three axioms, then a ground, G0, that is stated and not derived; every later section is a theorem
(T) or a derived constraint (C) that names the axioms it rests on~\cite{dmp}. We state the parts this paper uses, quoting DMP's sentences
where the wording matters, and give each a formal counterpart \tagD.

\subsection{Words}

The undefined words are \emph{part}, \emph{act}, \emph{witness}, \emph{evidence} and \emph{time}. A part is ``a Duality of Trust: two halves,
one owning trust and one owning information; nothing further is said about what a part is made of''. Trust is a relation $T(w,p,a,t)$:
witness $w$ trusts part $p$ for act $a$ at time $t$. Two witnesses are \emph{independent} when they share no member of the failure classes
declared for their layer (host, key, author, engine, person, time source, and whatever else the minting lists). A \emph{layer} is a place
where trust is decided; its two witnesses are its \emph{halves}, the \emph{trust half} (whether an act may happen) and the
\emph{information half} (whether it is meant and right). A \emph{record} is evidence that outlives its act. A boundary \emph{contains}
information when no part beyond it is trusted for any act on that information, so that no act carries it across.

\begin{definition}[gate]\label{def:gate}
Let $\mathcal V=\{\textsf{refuse}<\textsf{hold}<\textsf{pass}\}$ be the verdict lattice with meet $\wedge$. A layer $L$ with halves $w_T$
(trust) and $w_I$ (information) decides an act $a$ at time $t$ by
\[
\operatorname{gate}_L(a,t)=v_{w_T}(a,t)\wedge v_{w_I}(a,t),
\]
and the act is performed exactly when $\operatorname{gate}_L(a,t)=\textsf{pass}$.
\end{definition}

\subsection{The three axioms}

\paragraph{Axiom 1 (the Duality).} A part acts only while the two halves of its layer, two independent witnesses, ``each with full reach
over the act and gated only by the other'', trust it for that act. ``One half owns trust and the other owns information; neither overrules
the other, and the two keep each other constantly in check and balance.'' Where they disagree the stricter answer stands; every input
gets exactly one answer. \emph{Formally}: Definition~\ref{def:gate}, with $w_T,w_I$ independent; the meet is ``the stricter answer
stands''. The axiom also carries the \emph{balance}: every part works towards the trust its boundary declares, restoring by evidence what
was lost without evidence against it.

\paragraph{Axiom 2 (Trust).} A witness trusts only on evidence it verified itself at the point of use; no information is ever evidence of
itself; the absence of evidence moves nothing. ``Trust is gravitational. It is a field, never a position: every standing has a source
(the witness that granted it), a depth (its grant chain back to the pair) and a freshness that decays.'' A boundary strong enough ``bends
trustspace so that no act carries information across it, and what still leaves is entropy and nothing else''. \emph{Formally}: each
verdict is a function $v_w(a,t)=J_w(E_w(a,t))$ of the evidence $E_w$ that $w$ itself verified at $t$, and a standing verified at $t_v$
with validity $\tau$ has freshness $\varphi(t)=\max\{0,1-(t-t_v)/\tau\}$.

\paragraph{Axiom 3 (Information).} ``Every grant, every use and every loss of trust leaves a record, written at the act, that every witness
of that act can read in full and that outlives the act. Any other reader sees the counts always and the content only by a grant. From trust
nothing can be derived about information but its existence.'' Content moves only by a grant to a part trusted for its class. Information
is hierarchical: every layer has its own Duality and is a tree rooted there; the layout is ``the median of that demand, recomputed by both
halves of the layer and never by one''. The axiom carries the mission ``while maintaining absolute Observability''. \emph{Formally}, the
existence clause is a noninterference condition~\cite{gm}: for every observable $O$ of the trust side and every two contents $x,x'$ that
agree on existence, $O(x)=O(x')$.

\subsection{The ground, containment and the field}

\paragraph{G0.} ``The dual gate cannot witness its own first instance: somebody has to trust the first two witnesses.'' That first act is
the \emph{minting}, in which at least five parts become the holders of the root keys; it is ``stated, not derived'', and every later
standing descends from it. Theorem T40: a minting originates only from G0 together with a state in which every axiom and theorem
provably holds.

\paragraph{Containment and the floor.} T37: information that no admission names is contained at the item itself; ``in the absence of any
external trust, trustspace collapses in on the item''. T31: contained information reaches no part by any act; the substrate still
radiates, and the design makes that emission minimal and \emph{thermal}, ``so it carries entropy and no information'' (the \emph{entropy
floor}); an outside observer learns of the inside ``only its size, its staleness and its rotation interval''.

\paragraph{The VIBE field.} T32 makes Axiom~2's field concrete. Each member $i$ of the system has a \emph{depth} $d_i$ (its shortest grant
chain back to the pair), a \emph{mass} $m_i$ (one plus the number of members deriving through it) and a \emph{shortfall} $s_i$ (what the
boundary declares for it minus what it holds, weighted by freshness). The VIBE field (Variable Influence Balancing Equation) and its
potential are
\begin{equation}\label{eq:vibe}
V_i=m_i\,s_i,\qquad \Phi=\sum_iV_i\ \ge 0,
\end{equation}
with $\Phi=0$ ``exactly when the boundary holds what it declares''. Under the axioms ``the system's own acts never raise the potential;
only losses and ageing do''. Every reader other than the two witnesses reads one number, the boundary's \emph{temperature}, its
mass-weighted staleness. We call $\Phi=0$ the \emph{balanced state}.

\subsection{Formal status}

A Lean~4 development, core library only, models the axioms from a byte-for-byte copy of DMP pinned by hash~\cite{lean4}. It holds 243
theorems on its main branch, each checked by Lean's kernel and by an independently written second kernel, and each resting on no axiom
beyond \texttt{propext}, \texttt{Quot.sound} and \texttt{Classical.choice}. Among them: the gate of a stack of layers is the meet along its
chain~\cite{denning}; every pass in a well-formed configuration was passed by a minted pair of independent witnesses, read as fail-prone
systems~\cite{mr}; adding a layer changes no existing decision (the frame rule~\cite{ohearn}); placement never enters a decision
(noninterference~\cite{gm}); the potential~\eqref{eq:vibe} never rises under the system's own refresh. The Lean concerns the axiom system,
not the physics: no theorem below rests on it, and it is cited here only so that the reader knows the axioms are consistent in at least one
model. Adversarial reviews of that development found, among others, that the axioms admit an inert model that performs no act, and that
independence is relative to a declared list of failure classes; neither finding touches what follows.

\section{The physical reading: two fields}\label{sec:fields}

\subsection{Hypotheses}

GUID reads DMP through five identifications. They are the only steps of the paper that no derivation supplies, and every result below
names the ones it uses. \nw{\vtag{} Two more follow in Section~\ref{sec:cosmo}, and version 5.1 refines the first and the sixth
(Hypotheses~1$'$ and~6$'$).}

\begin{hypothesis}[trust is existence]\label{h:exist}
DMP's trust half is \emph{existence}: the field of Axiom~2 is the gravitational field, and a standing in it is the existence of what it
stands for. \tagH
\end{hypothesis}

\begin{hypothesis}[existence is measured by energy-momentum]\label{h:energy}
The existence of a body is measured by its energy-momentum tensor $T_{\mu\nu}$: two bodies with the same $T_{\mu\nu}$ have the same
existence, whatever they are made of. \tagH
\end{hypothesis}

\begin{hypothesis}[the existence field is a frame]\label{h:frame}
The existence field spans every direction of spacetime, three of space and one of time: it is a frame of four vector fields
$e_a{}^{\mu}$ (a tetrad), not a single vector. \tagH
\end{hypothesis}

\begin{hypothesis}[information carries quantum content]\label{h:info}
DMP's information half carries the quantum content of the world, the states of matter and radiation, together with one scalar field
$\psi$ that sets the strength of the electromagnetic coupling, $\alpha=\alpha(\psi)$. \tagH
\end{hypothesis}

\begin{hypothesis}[containment is a horizon]\label{h:horizon}
A containing boundary is an event horizon; the counts every outside reader sees are its asymptotic charges (mass, angular momentum,
electric charge); the entropy floor is the horizon's thermal radiation. \tagH
\end{hypothesis}

\begin{vnew}
\begin{hypothesisprime}[the gravitational plane; refines Hypothesis~\ref{h:exist}]
\vtag{} The trust field of Axiom~2 (``Trust is gravitational. It is a field'') is the \emph{gravitational plane}: the frame of
Hypothesis~\ref{h:frame}. It is separate from the existence and the information it contains and is shaped by them only through the counts
every reader sees (Axiom~3), their energy and momentum; it is never either of them. Its link to G0 is the causal chain of whatever exists.
Existence is what is granted, the rest mass the Higgs vacuum gives (Section~\ref{sec:higgsplit}); information is the content. \tagH
\end{hypothesisprime}

\noindent\vtag{} \emph{Why the refinement.} Hypothesis~\ref{h:exist} made gravity the existence half itself. The data already pointed
elsewhere: free fall weighs all energy, granted and bound alike (Section~\ref{sec:fix}, \chk{Q6}), and a gravity that were existence alone
would weigh only the granted share. With Hypothesis~1$'$ every result of Section~\ref{sec:fix} keeps its proof, the plane taking the place
of the existence frame, and two proofs become shorter (the remarks after Proposition~\ref{pr:wep} and Lemma~\ref{lem:classical}). The
refinement keeps the classical-gravity prediction only because the plane is a DMP structure, shaped by the counts; were it outside DMP's
structures altogether, DMP would say nothing about whether gravity is quantum, and the test of Section~\ref{sec:exp} would be lost.
\end{vnew}


\subsection{The existence frame}

A tetrad is a set of four one-forms $e^a=e^a{}_\mu\,\dd x^\mu$, $a=0,\dots,3$, orthonormal with respect to the metric they define:
\begin{equation}\label{eq:tetrad}
g_{\mu\nu}=\eta_{ab}\,e^a{}_\mu\,e^b{}_\nu,\qquad \eta=\operatorname{diag}(-1,1,1,1).
\end{equation}
At each point the tetrad is fixed up to a local Lorentz transformation $e^a\to\Lambda^a{}_b(x)\,e^b$, which leaves $g_{\mu\nu}$ unchanged.
The tetrad has 16 components; the local Lorentz group removes 6; the 10 that remain are the components of the metric. General relativity
can be written with the tetrad as its only field, either in the Einstein--Cartan form with a Lorentz connection or in the teleparallel
form, with exactly Einstein's predictions~\cite{aldrovandi} \tagP. Hypothesis~\ref{h:frame} matters: a single timelike unit vector field
added to the metric would single out a rest frame, which is Einstein--aether theory~\cite{aether}; its preferred-frame effects are bounded
by the simultaneous arrival of gravitational waves and light from GW170817 to $|c_{\text{gw}}/c-1|\lesssim10^{-15}$~\cite{gw170817}
\chk{Q17}. A frame of four has no preferred direction.

\subsection{The information field and the content}

The information half holds every quantum field of matter and radiation, which this paper calls \emph{content}, and one scalar $\psi$ whose
value sets the fine-structure constant. A scalar that multiplies the Maxwell term and leaves the gravitational term alone is Bekenstein's
theory of a varying $\alpha$~\cite{bekenstein82}; with the scalar's energy fed back into cosmology it is the theory of Sandvik, Barrow
and Magueijo~\cite{bsbm} \tagP. The action whose shape the hypotheses fix is
\begin{equation}\label{eq:action}
S=\int\dd^4x\,|e|\left[\frac{c^4}{16\pi G}R(e)-\frac{c^4}{16\pi G}\,\partial_\mu\phi\,\partial^\mu\phi-\frac{1}{4\mu_0}\,B(\phi)\,
F_{\mu\nu}F^{\mu\nu}\right]+S_{\text{content}}[e,\phi,\Psi],
\end{equation}
where $|e|=\det e^a{}_\mu=\sqrt{-g}$, $R(e)$ is the curvature scalar of the metric~\eqref{eq:tetrad}, $\phi$ is $\psi$ measured in units of
the reduced Planck mass (so that its kinetic term is normalised like the gravitational one), $\alpha\propto1/B(\phi)$, and $\Psi$ stands for
the quantum fields of matter. We write
\begin{equation}\label{eq:lambda}
\lambda\equiv\frac{\dd\ln\alpha}{\dd\phi}
\end{equation}
for the scalar's coupling. The action is our construction; what the axioms fix about it is the subject of the next section.

\subsection{The dictionary}

\begin{center}\small
\begin{tabularx}{\linewidth}{@{}lLL@{}}
\toprule
DMP & GUID (physics) & Hypothesis \\
\midrule
trust half; the trust field (Axiom 2) & the existence frame $e^a{}_\mu$; gravity & \ref{h:exist}, \ref{h:frame} \\
information half & quantum content $\Psi$ and the scalar $\phi$ that sets $\alpha$ & \ref{h:info} \\
a standing; mass of a member (T32) & existence, measured by $T_{\mu\nu}$ & \ref{h:energy} \\
record, read in full by every witness (Axiom 3) & the value of the frame, held by every system it acts on & \ref{h:exist} \\
counts, seen by every reader (Axiom 3) & asymptotic charges $M$, $J$, $Q$ & \ref{h:horizon} \\
containing boundary; trap (T27, T31) & event horizon & \ref{h:horizon} \\
entropy floor; temperature of a boundary (T31, T32) & thermal radiation of the horizon; Hawking temperature & \ref{h:horizon} \\
G0, the minting & the start: inflation with the decay of its oscillons; the classical era begins as it completes & \ref{h:minting} \\
balanced state $\Phi=0$ (T32) & the vacuum at the minimum of its potential & Section~\ref{sec:matter} \\
\nw{\vtag{} trust field (Axiom 2)} & \nw{the gravitational plane (the frame), separate from what it contains and shaped by its counts} &
\nw{1$'$} \\
\nw{\vtag{} existence} & \nw{the rest mass the Higgs vacuum grants} & \nw{1$'$; Section~\ref{sec:higgsplit}} \\
\bottomrule
\end{tabularx}
\end{center}

\section{What the axioms fix}\label{sec:fix}

\subsection{Free fall is universal}

\begin{proposition}[weak equivalence principle]\label{pr:wep}
Under Hypotheses~\ref{h:exist} and~\ref{h:energy}, two bodies with the same energy-momentum fall alike in the existence frame, whatever
they are made of.
\end{proposition}

\begin{proof}
Axiom~3 \tagD: from trust nothing can be derived about information but its existence. By Hypothesis~\ref{h:exist} the existence frame is
the trust half, so no response of the frame to a body can depend on what the body carries (composition, internal state) except through
the body's existence. By Hypothesis~\ref{h:energy} existence is $T_{\mu\nu}$. So the frame couples to a body through $T_{\mu\nu}$ alone,
and two bodies equal in it follow the same worldline: the weak equivalence principle \tagP.
\end{proof}

\begin{vnew}
\noindent\vtag{} \emph{Under Hypothesis~1$'$} the proof needs no reading of existence: the plane is shaped only by the counts every reader
sees (Axiom~3), which are energy and momentum whatever carries them, so free fall is universal, and the plane weighs granted and bound
energy alike, as MICROSCOPE requires.
\end{vnew}


The principle is among the best-tested in physics. The MICROSCOPE satellite compared the free fall of titanium and platinum and found the
E\"otv\"os ratio $\eta=[-1.5\pm2.3\,(\text{stat})\pm1.5\,(\text{syst})]\times10^{-15}$~\cite{microscope}. Two consequences are sharp.

\emph{Gravity weighs all energy.} About 9\% of a nucleon's mass comes from the Higgs-given quark masses; the rest is the energy of quark
motion and of the gluon field~\cite{yang}. A gravity that weighed only the Higgs-given mass would make titanium and platinum fall
differently by $\eta\approx1\times10^{-3}$ \chk{Q6}, more than $10^{11}$ times what MICROSCOPE allows. The frame weighs the binding in full.

\emph{The information scalar may barely pull.} A scalar that sets $\alpha$ changes the electrostatic energy of every nucleus, so each body
of mass $m_A$ carries a scalar charge $\partial\ln m_A/\partial\phi=\lambda f_A$, where $f_A$ is the electrostatic share of its mass. With the
normalisation of~\eqref{eq:action} the exchange of $\phi$ adds to the Newtonian force between $A$ and $B$ a fraction
\begin{equation}\label{eq:fifth}
\frac{F_\phi}{F_N}=2\lambda^2f_Af_B ,
\end{equation}
so two bodies falling towards the Earth (share $f_\oplus$) differ by $\eta=2\lambda^2f_\oplus(f_1-f_2)$ (Dvali and Zaldarriaga~\cite{dz}).
With the Coulomb term of the semi-empirical mass formula, $f_{\text{Ti}}=2.0\times10^{-3}$ and $f_{\text{Pt}}=4.1\times10^{-3}$, and
MICROSCOPE gives
\begin{equation}\label{eq:lambdabound}
|\lambda|<2\times10^{-5}\qquad\chk{Q6}.
\end{equation}
\emph{Verdict: derived} (from Axiom~3 with Hypotheses~\ref{h:exist},~\ref{h:energy}); \emph{holds} against MICROSCOPE.

\subsection{The scalar stays out of gravity: \texorpdfstring{$\gamma=1$}{gamma = 1}}

\begin{proposition}[light bending]\label{pr:gamma}
If the information half cannot change how the existence frame responds to existence, the parametrised post-Newtonian parameter is
$\gamma=1$ exactly, as in general relativity.
\end{proposition}

\begin{proof}
Axiom~1 \tagD: neither half overrules the other. A scalar that multiplies the curvature term, $F(\phi)R$, sets the local strength of gravity
and so overrules the existence half; Axiom~1 excludes it. What remains is minimal coupling, as in~\eqref{eq:action}: the scalar's energy is
one more source of the frame, negligible in the solar system, and the post-Newtonian expansion is that of general relativity,
$\gamma=1$~\cite{will} \tagP.
\end{proof}

In Brans--Dicke gravity~\cite{bd}, where the scalar does multiply the curvature, $\gamma=(1+\omega)/(2+\omega)$, and the Cassini
measurement $\gamma-1=(2.1\pm2.3)\times10^{-5}$~\cite{cassini} forces $\omega>4\times10^4$ at two standard deviations \chk{Q7}: the
scalar must decouple almost completely. Axiom~1 asks for that decoupling from the start. \emph{Verdict: derived}; \emph{holds}.

\subsection{The frame is classical}

We need two theorems of quantum information.

\begin{definition}[broadcasting]
A family of states $\{\rho_s\}$ on a system $S$ can be \emph{broadcast} if there is one physical process (a completely positive,
trace-preserving map) $\mathcal B:S\to S_1\otimes S_2$ such that both reduced states of $\mathcal B(\rho_s)$ equal $\rho_s$, for every $s$.
For pure states and a product output this is \emph{cloning}.
\end{definition}

\begin{lemma}[no-cloning; no-broadcasting]\label{lem:nobroadcast}
(i) A unitary that copies two pure states $\ket s,\ket t$, $U\ket s\ket0=\ket s\ket s$ and $U\ket t\ket0=\ket t\ket t$, exists only if
$\langle s|t\rangle\in\{0,1\}$ (Wootters and Zurek~\cite{wz}; Dieks~\cite{dieks}). (ii) A family of states can be broadcast if and only if
its members commute pairwise (Barnum, Caves, Fuchs, Jozsa and Schumacher~\cite{barnum}). \tagP
\end{lemma}

\begin{proof}[Proof of (i)]
Unitaries preserve inner products, so $\langle s|t\rangle=\langle s|t\rangle^2$, whose only solutions are $0$ and $1$. For (ii) see~\cite{barnum};
the ``if'' direction is the copier in the common eigenbasis, $\ket i\to\ket i\ket i$, which \chk{Q3} runs: a diagonal state is broadcast,
$\ket+$ is not (each reader is left with $I/2$).
\end{proof}

\begin{lemma}[the existence frame is classical]\label{lem:classical}
Under Hypothesis~\ref{h:exist}, every value the existence frame takes after the first recorded act belongs to a pairwise commuting family:
the frame admits no superposition of two of its values.
\end{lemma}

\begin{proof}
Axiom~3 \tagD: every grant, use and loss of trust leaves a record that every witness of the act can read in full, and any other reader sees
the counts always. Axiom~2 \tagD: every standing of the trust field is evidence held in such records. Axiom~1 \tagD: there are at least two
independent witnesses. So every value of the trust field is held in full by at least two independent readers: it is broadcast. By
Lemma~\ref{lem:nobroadcast}(ii) the values commute. By Hypothesis~\ref{h:exist} they are the values of the existence frame.
\end{proof}

\begin{remark}
The lemma is close to definitional: DMP demands absolute observability of its shared structure, and the no-broadcasting theorem states the
price. It is not a discovery about nature; it is the statement of which side of an open question the axioms take. It has a precise physical
counterpart in quantum Darwinism: a quantity of which the environment holds many redundant records is objective and classical, and generic
quantum dynamics makes exactly the pointer observables so recorded~\cite{zurek,bph}. DMP's ``every witness reads in full'' is that
redundancy, imposed as an axiom. Read in the other direction, the lemma is a design statement for information systems: a control structure
that every party can audit in full is classical, while the data it governs need not be.
\end{remark}

\begin{vnew}
\noindent\vtag{} \emph{Under Hypothesis~1$'$} the lemma applies to the plane directly: the plane depends on the parts only through the
counts, which every reader holds in full, so from the parts' point of view it is a measure-and-prepare channel (Lemma~\ref{lem:locc}) and
classical.
\end{vnew}


\subsection{A classical mediator cannot entangle}

\begin{lemma}[classical mediation is LOCC]\label{lem:locc}
Let $A$ and $B$ be quantum systems that interact only through a mediator $M$, by alternating local interactions $A$--$M$ and $M$--$B$. If
$M$ is classical (its states commute pairwise in a fixed basis $\{\ket k\}$ at every step), then the induced map on $AB$ is a local
operation with classical communication, and it maps every separable state of $AB$ to a separable state.
\end{lemma}

\begin{proof}
A classical $M$ is read in its basis after each interaction. An $A$--$M$ interaction followed by that reading is an instrument on $A$: it
takes $\rho_A$ to the ensemble $\{p_k,\rho_A^{(k)}\}$ with the outcome $k$ written in $M$. An $M$--$B$ interaction with $M$ in $\ket k$ is a
channel $\mathcal F_k$ on $B$. One round therefore maps $\rho_{AB}$ to $\sum_k(\mathcal E_k\otimes\mathcal F_k)(\rho_{AB})$ with
$\mathcal E_k$ completely positive and $\sum_k\mathcal E_k$ trace preserving: a one-way LOCC map. Further rounds compose LOCC maps, which
remain LOCC. LOCC maps take separable states to separable states, and no entanglement measure increases under them~\cite{vw} \tagP.
\end{proof}

Galley, Giacomini and Selby prove the same statement in the generality of probabilistic theories: a mediator that entangles two quantum
systems by local interaction cannot be classical~\cite{ggs}. Check \chk{Q4} runs Lemma~\ref{lem:locc} on qubits: three random local
couplings $A$--$M$, $M$--$B$, $A$--$M$ leave $A$ and $B$ entangled in most of 300 trials when $M$ is coherent (median negativity $0.12$)
and never when $M$ is dephased after each coupling (largest negativity below $10^{-12}$).

\subsection{DMP hosts the proposed theory}

\begin{definition}[base configuration, host]\label{def:model}
The \emph{base configuration currently assumed for our universe} is what physics supplies and DMP does not derive, the configuration in
DMP's sense (T4): the constants $G$, $\hbar$, $c$, $\Lambda$ and $\alpha$, the fields and parameters of the Standard Model, and the initial
data of the $\Lambda$CDM model. A configuration that obeys DMP \emph{hosts} a theory of gravity if its own shared structures (the trust
field, the layout, the order of acts, the counts) carry the theory's classical degrees of freedom and its information half carries the
theory's quantum content, so that the configuration reproduces everything the theory predicts.
\end{definition}

\begin{theorem}[DMP hosts the proposed theory of quantum gravity]\label{th:host}
Under Hypotheses~\ref{h:exist}--\ref{h:info}, and supplied with the base configuration currently assumed for our universe, DMP hosts the
theory of this paper: a classical existence frame over quantum content. In it, (i) the frame, carried by DMP's shared structures, mediates
no entanglement between the systems it acts on; (ii) the information half carries any quantum content.
\end{theorem}

\begin{proof}
(i) DMP's shared structures are broadcast records (Axiom~3), so by Lemma~\ref{lem:classical} the field they carry is classical, whichever
structure carries it, and by Lemma~\ref{lem:locc} a classical mediator creates no entanglement. (ii) Axiom~3 restricts what is
\emph{derived from trust} and how content moves (by grant); it does not require content to be broadcast, so nothing in the axioms forbids
non-commuting states in the information half \tagD. A grant moves content without copying it, as no-cloning requires. The base
configuration supplies the constants and the content that DMP does not derive.
\end{proof}

Two remarks place the theorem. First, the constants: DMP derives none of them, and the theorem needs them supplied, as every physical
theory needs its parameters (Section~\ref{sec:break}). Second, universality: whatever can run a universal quantum computer as content can
also run a theory in which the geometry itself is quantum, by the Church--Turing--Deutsch principle~\cite{deutsch} (Section~\ref{sec:host});
in such a theory the quantum geometry would sit in the information half and would not be DMP's frame. Part~(i) is therefore the proposed
theory's characteristic, falsifiable prediction, and Section~\ref{sec:exp} names the experiments that test it.

\subsection{Gravitational waves and the graviton}

The frame's linearised excitations, $g_{\mu\nu}=\eta_{\mu\nu}+h_{\mu\nu}$, are a massless spin-2 field. In $D$ spacetime dimensions such a
field has $D(D-3)/2$ propagating polarisations, two in 3+1 \chk{Q8}. Weinberg's low-energy theorem~\cite{weinberg64} and Deser's
self-coupling argument~\cite{deser70} show that a massless spin-2 field must couple equally to all energy and, made consistent, is general
relativity \tagP: the spin-2 wave is the existence frame's own, and its universal coupling is Proposition~\ref{pr:wep} again.

By Lemma~\ref{lem:classical}, once the start is complete (Section~\ref{sec:cosmo}) these waves are classical waves of the frame: GUID has no graviton as a quantum
of the frame \tagD. That is consistent with every observation, for reasons the literature has made precise. Dyson showed that detectors of
the LIGO type cannot register single gravitons, because mirrors heavy enough would collapse into a black hole~\cite{dyson}; Rothman and
Boughn found that even a Jupiter-mass absorber near a neutron star would be swamped by neutrinos, and a shield against them would collapse
too~\cite{rb}. Every gravitational wave measured is deep in the classical regime: the energy GW150914 radiated corresponds to more than $10^{78}$
quanta at its frequency \chk{Q9}. Localising any quantum inside its own Schwarzschild radius needs a wavelength below $\sqrt{4\pi}\,L_P$ and
an energy above $\sqrt\pi\,E_P$ \chk{Q9}, the Planck scale at which measurement and gravity cannot be separated~\cite{mead}; a LIGO-band
graviton is $2\times10^{-41}$ of that. The proposals to register single quanta of a gravitational wave as jumps of a cooled resonant
bar~\cite{tobar} are a gravitational photoelectric effect, and like the photoelectric effect they can be explained by a classical wave acting
on a quantised detector~\cite{cdr}. So the absence of the graviton in GUID is not testable by counting gravitons; it is testable only
through the mediation experiments of Section~\ref{sec:exp}. \emph{Verdict: derived}; \emph{consistent} with all data; testable indirectly.

\section{How a classical frame can carry quantum matter}\label{sec:cq}

Theorem~\ref{th:host} places GUID among the theories in which spacetime is classical and matter quantum. Such theories have a known
difficulty, and its resolution is a prediction.

\subsection{The semiclassical equation fails}

The simplest coupling, $G_{\mu\nu}=(8\pi G/c^4)\langle\hat T_{\mu\nu}\rangle$ (M{\o}ller~\cite{moller}, Rosenfeld~\cite{rosenfeld63}),
sources gravity by the quantum expectation value. A mass in an equal superposition of two places then attracts a probe towards the midpoint,
where the mass never is. Page and Geilker tested a macroscopic version, with the position of lead masses set by a quantum random event, and
found the probe following the actual masses, not the average~\cite{pagegeilker} \tagP. The equation is also nonlinear in the quantum state,
which in general permits signalling. So a classical frame cannot simply be sourced by expectation values.

\subsection{Consistent classical-quantum dynamics}

The consistent alternative is a joint state $\varrho(z)$ of a classical phase-space variable $z$ (here the frame) and a quantum system,
evolving by a map that is linear in $\varrho$, completely positive, and preserves normalisation, as probability theory demands. Oppenheim
showed that such dynamics exists and can couple a classical metric to quantum fields~\cite{oppenheim}. For continuous, Markovian
classical-quantum dynamics the generator has, besides the usual Hamiltonian part, three kinds of term: a \emph{decoherence} term on the
quantum side with coefficient matrix $D_0$, a \emph{back-reaction} (drift) of the classical variable driven by the quantum state with
coefficients $D_1$, and a \emph{diffusion} of the classical variable with coefficients $D_2$. Complete positivity forces the
trade-off~\cite[Eq.~(23)]{osswd}
\begin{equation}\label{eq:tradeoff}
D_1\,D_0^{-1}\,D_1^{\dagger}\;\preceq\;2D_2 ,
\end{equation}
with $D_0^{-1}$ the generalised inverse, and in addition $D_0$ cannot vanish where $D_1$ does not \tagP. In words: if quantum matter pushes on
the classical frame at all ($D_1\neq0$), the matter must decohere ($D_0>0$) and the frame must diffuse ($D_2>0$), and the product of the two
is bounded below by the square of the coupling. Long coherence times of massive superpositions demand strong random fluctuations of the
gravitational field, and weak fluctuations demand short coherence times.

\begin{corollary}[GUID diffuses]\label{co:diffuse}
If the existence frame (Lemma~\ref{lem:classical}) is sourced by quantum content (Hypothesis~\ref{h:info}) through a dynamics that is linear,
completely positive and Markovian, then mass superpositions decohere at a non-zero rate and the gravitational field fluctuates randomly,
with~\eqref{eq:tradeoff} bounding the product from below.
\end{corollary}

The corollary turns the classicality of the frame from a philosophical position into a quantitative one: in the Newtonian limit, the
decoherence rates of massive superpositions measured in interferometry and the noise in precision measurements of gravitational forces
together squeeze the remaining parameter space from both sides~\cite{osswd}.

\subsection{The record is a measurement, the layout is feedback}

There is a second, older family of consistent classical-quantum models, and DMP's architecture coincides with it. Kafri, Taylor and Milburn
model the gravitational interaction between two masses as a continuous weak measurement of each mass's position followed by a classical
force applied to the other according to the recorded result~\cite{ktm}; Tilloy and Di\'osi source a classical Newtonian potential from a
continuously measured (spontaneously collapsing) mass density~\cite{tilloydiosi}. In both, gravity reproduces Newton's law on average,
cannot entangle, decoheres superpositions and adds noise, as Corollary~\ref{co:diffuse} requires.

Read through the dictionary, DMP's Axiom~3 has exactly this structure \tagD+\tagH. Every act leaves a record that every witness reads in
full: a measurement of the act's existence. Any other reader sees the counts: the measured mass density, coarse-grained. The layout of the
system is ``the median of that demand, recomputed by both halves of the layer and never by one'': a classical function of the records,
applied back to the system, which is feedback. From the parts' point of view the whole loop is a measure-and-prepare channel, which is why
it cannot entangle (Lemma~\ref{lem:locc}).

\begin{proposition}[DMP is a measurement-and-feedback gravity]\label{pr:ktm}
Under Hypotheses~\ref{h:exist}--\ref{h:info}, the dynamics that DMP's axioms prescribe for the trust field is a continuous measurement of
existence (energy-momentum) followed by classical feedback on the geometry, of the Kafri--Taylor--Milburn and Tilloy--Di\'osi type.
\end{proposition}

The proposition fixes the architecture, not the constants. A measurement-and-feedback model has a free parameter, the measurement strength
or equivalently the smearing length of the measured mass density, which sets the decoherence rate and the noise. DMP derives no constant
(Section~\ref{sec:break}), and the most economical choice, the parameter-free Di\'osi--Penrose collapse~\cite{diosi87,penrose96}, is
excluded by the absence of the spontaneous radiation it predicts in an underground germanium detector~\cite{donadi} \tagP. GUID must
therefore carry a free strength, fixed by experiment or by a principle that ties DMP to physical units. \emph{Verdict: derived} as a
structure; Section~\ref{sec:bridges} computes the constant from the two principles GUID suggests.

\subsection{Fixing the free constant}\label{sec:bridges}

Proposition~\ref{pr:ktm} fixes the architecture, not the constant. DMP cannot supply the constant by itself: its text makes every
parameter configuration (T4, T32), and it has no unit of length, time or action. A number therefore needs one bridge principle that ties a
DMP structure to a physical scale. We compute the two that GUID's own dictionary suggests. Neither is derived from the axioms; both are
hypotheses \tagH, stated so that they can fail.

\paragraph{The balance rule.} Both bridges use one structural rule. Axiom~1's balance moves every part towards its declared state, and the
system's own acts never raise the VIBE potential (T32); read as dynamics, the frame adds no more noise than consistency requires. GUID
therefore holds the trade-off~\eqref{eq:tradeoff} at equality,
\begin{equation}\label{eq:saturate}
2D_2=D_1\,D_0^{-1}\,D_1^{\dagger},
\end{equation}
so that a measured gravitational decoherence would fix the spacetime diffusion, and conversely. Other classical-quantum theories leave the
two free above the bound; \eqref{eq:saturate} is GUID's own rule.

\paragraph{Bridge A: the outermost horizon sets the scale.} The stochastic remnant of Section~\ref{sec:darkmatter} depends on the outer
cut-off $r_{\max}$ of the path integral. Oppenheim and Russo find, in units with $c=1$, the mean
$\gamma_1=-\tfrac92\,\gamma_2\,r_{\max}\,(1-4\beta)/(5-18\beta)$, with $\gamma_2=\Lambda/3$ and a parameter $\beta<0$ of their action to which
the result is insensitive~\cite{oppenheimrusso}. In GUID the outermost containing boundary is the cosmological event horizon, which in
$\Lambda$CDM tends to the de Sitter radius $r_{\text{dS}}=c/H_\infty=1.66\times10^{26}$\,m. With $r_{\max}=r_{\text{dS}}$ and $c$ restored,
\begin{equation}\label{eq:bridgeA}
|\gamma_1|=(0.90\ \text{to}\ 1.00)\;cH_\infty=(4.9\ \text{to}\ 5.4)\times10^{-10}\,\text{m/s}^2,
\end{equation}
with no free parameter \chk{Q19}. It has the right order of magnitude and is too strong. Matching the turnover speed~\eqref{eq:btf} to the
observed baryonic Tully--Fisher relation $v^4=GMa_0$ requires $|\gamma_1|=a_0/2=6\times10^{-11}$\,m/s$^2$; Bridge~A exceeds this by a factor
of 8 in $v^4$, 1.7 in $v$. For a galaxy of $10^{11}$ solar masses it gives 340\,km/s at $r_*=7.6$\,kpc, where the relation gives 200\,km/s,
and the linear term lets the curve rise to 620\,km/s at 50\,kpc, where observed curves stay flat \chk{Q19}. \emph{Verdict: refuted} as it
stands. Its near miss is Milgrom's coincidence $a_0\approx cH/2\pi$, which is known.

\paragraph{Bridge B: one record per Planck time.} The second bridge reads DMP's tick as the Planck time: the frame records existence at
Planck strength, so the coupling takes its value in Planck units, and the balance rule~\eqref{eq:saturate} holds. In the Newtonian limit
Oppenheim et al.\ bound the coupling from both sides, from above by precision measurements of gravitational forces and from below by the
observed coherence of massive molecules~\cite{osswd}. At Planck strength the discrete version of the dynamics has
$(l_P^3/m_P)D_2=m_Pt_P\approx10^{-51}$\,kg\,s, below the allowed window $10^{-25}$ to $10^{-1}$\,kg\,s: \emph{excluded}. The continuous
version with the Di\'osi--Penrose kernel has $l_P^2D_2=m_P^2/c=1.6\times10^{-24}$\,kg$^2$\,s/m, inside the allowed window $10^{-40}$ to
$10^{-9}$\,kg$^2$\,s/m: \emph{allowed} \chk{Q19}. With the trade-off at equality, a mass $M$ of size $R$ in a superposition of two
well-separated positions decoheres at the rate
\begin{equation}\label{eq:dprate}
\lambda=\frac{M^2}{R\,l_P^2D_2}=\frac{GM^2}{\hbar R},
\end{equation}
since $c/m_P^2=G/\hbar$: exactly the Di\'osi--Penrose rate~\cite{diosi87,penrose96}, to within the order-one factors of the Newtonian
estimate. For a fullerene it is $6\times10^{-16}$ per second, far below what interferometry allows. For the diamond microsphere of the Bose
et al.\ proposal ($10^{-14}$\,kg, radius 0.88\,\textmu m) it is 72 per second: coherence is lost in about 14\,ms, while the entangling phase
needs 8 to 25\,s, so only $2\times10^{-3}$\,rad builds up \chk{Q19}. One length remains that the bridge does not fix: the mass density of each
nucleus must be smeared over at least about $10^{-10}$\,m, or the noise makes matter radiate more than an underground germanium detector
sees~\cite{donadi}, and later analyses quote $4\times10^{-10}$\,m~\cite{figurato}; at such smearing the bulk term in~\eqref{eq:dprate} still
dominates \chk{Q19}. \emph{Verdict: live} and sharp: with Bridge~B, mesoscopic superpositions decohere at the Di\'osi--Penrose rate, and a
Bose--Marletto--Vedral experiment sees decoherence and no entanglement.

\paragraph{What the bridges show.} Made parameter-free with its own principles, GUID lands on known physics: Bridge~A is refuted, and
Bridge~B is the Di\'osi--Penrose model, which predates this work. What GUID adds is a reason for that model, the balance
rule~\eqref{eq:saturate} together with a record at Planck strength. A number of its own would need a prediction in which GUID differs from
both the classical-quantum theory and Di\'osi--Penrose. The most promising lead is DMP's two independent witnesses, which may imply two
independent noise sources with a statistical signature that a single source lacks; we have not worked it out.

\section{The experiments that decide}\label{sec:exp}

\subsection{Gravitationally mediated entanglement}

Bose et al.~\cite{bose} and Marletto and Vedral~\cite{mv} proposed to place two masses, each in a superposition of two positions, close
enough that only gravity couples them, and to test whether they end entangled. Let each mass $m$ be in $(\ket L+\ket R)/\sqrt2$ with the two
positions $\Delta x$ apart and the centres $d$ apart on one line. A coherent Newtonian field gives each branch $(i,j)$ the phase
$\varphi_{ij}=Gm^2t/(\hbar d_{ij})$, where $d_{ij}$ is the distance between the masses in that branch, so after time $t$
\begin{equation}\label{eq:bmv}
\ket{\Psi(t)}=\tfrac12\sum_{i,j\in\{L,R\}}e^{i\varphi_{ij}(t)}\ket i_A\ket j_B .
\end{equation}
Up to local phases this is $\tfrac12(\ket{00}+\ket{01}+\ket{10}+e^{i\Delta\varphi}\ket{11})$ with
$\Delta\varphi=\varphi_{RL}+\varphi_{LR}-\varphi_{LL}-\varphi_{RR}$; its concurrence is $|\sin(\Delta\varphi/2)|$ and its negativity
\begin{equation}\label{eq:neg}
\mathcal N=\tfrac12\,|\sin(\Delta\varphi/2)|.
\end{equation}
For the parameters of~\cite{bose}, $m=10^{-14}$\,kg, $\Delta x=250$\,\textmu m, $d=450$\,\textmu m, the entangling phase grows at
$0.126$\,rad/s: one radian after $8.0$\,s, maximal entanglement ($\mathcal N=1/2$, $\Delta\varphi=\pi$) after $25$\,s \chk{Q5}. A recorded,
classical field (Lemma~\ref{lem:classical}) leaves each branch's phase but removes the coherence between branches, and the masses stay
separable at all times, $\mathcal N=0$ \chk{Q5}.

\textbf{GUID predicts $\mathcal N=0$.} With Bridge~B of Section~\ref{sec:bridges} it predicts more: the masses decohere at the
Di\'osi--Penrose rate, about 72 per second for the proposed microspheres, long before an entangling phase can build up. A positive result would show that the gravitational field is not the existence frame: the identification
of Hypothesis~\ref{h:exist} would fail for nature and the proposed theory would be refuted; a quantum geometry would then run on DMP only
as content (Theorem~\ref{th:host}(ii)). Two caveats are
under debate. Aziz and Howl argue that classical gravity coupled to quantum matter fields can produce entanglement through virtual matter
processes, at a rate far below the quantum-gravity rate~\cite{azizhowl}; Marletto, Oppenheim, Vedral and Wilson~\cite{movw},
Di\'osi~\cite{diosi25}, Schneider, Huggett and Linnemann~\cite{shl} and Gundhi, Infantino and Bassi~\cite{gib} dispute it. And
Christodoulou et al.\ showed that the entanglement follows from linearised quantum gravity on a fixed background with local
mediation~\cite{christodoulou23}, which is the content-side mechanism of Section~\ref{sec:holo}'s dissent. Neither changes the prediction:
in GUID the frame is recorded, and Lemma~\ref{lem:locc} applies to it.

\subsection{Quantumness without entanglement}

Lami, Pedernales and Plenio showed that the quantum nature of gravity can be tested without creating entanglement: an LOCC inequality bounds
how well any local classical channel can reproduce a gravitational evolution of coherent states of massive oscillators, such as torsion
pendula, and a measured violation would exclude a classical mediator~\cite{lpp}. The test needs no large spatial superpositions. GUID
predicts that the inequality holds.

\subsection{Diffusion against decoherence}

Corollary~\ref{co:diffuse} gives a two-sided test. Interferometry with ever heavier molecules and levitated particles bounds the decoherence
that gravity may cause from above; precision measurements of the gravitational force, which has been measured between masses as small as
90\,mg~\cite{westphal}, bound the random fluctuations of the field from above. Inequality~\eqref{eq:tradeoff} says the two cannot both be
made arbitrarily small. GUID predicts that they cannot.

\subsection{What existing experiments already show}

Quantum matter responds to gravity as the framework requires: neutrons in a gravitational potential acquire the phase $mgA/\hbar v$
predicted for quantum matter in a classical field (Colella, Overhauser and Werner~\cite{cow}), and ultracold neutrons above a mirror occupy
discrete gravitational quantum states~\cite{nesvizhevsky}. In both, the matter is quantum and the field is a given potential; neither tests
whether the field itself can be in superposition. The current state of the field is reviewed by Carney, Stamp and Taylor~\cite{cst}.
\emph{Verdict on GUID's classical frame: live, falsifiable.}

\section{The thermodynamics of the frame}\label{sec:thermo}

\subsection{The existence of a boundary is its temperature}

GUID defines the existence a boundary asserts as the product of the local strength of gravity $g$, Planck's constant $h$ and the
inverse of the speed of light, which DMP's reading takes as the universal tick rate:
\begin{equation}\label{eq:Eex}
E_{\text{ex}}\equiv\frac{g\,h}{c}.
\end{equation}

\begin{proposition}[existence and the Unruh temperature]\label{pr:unruh}
$E_{\text{ex}}$ is the only energy that an acceleration, Planck's constant and the speed of light can form, and
\begin{equation}\label{eq:unruh}
E_{\text{ex}}=4\pi^2\,k_BT_U,\qquad k_BT_U=\frac{\hbar g}{2\pi c},
\end{equation}
where $T_U$ is the temperature an observer of proper acceleration $g$ assigns to the vacuum (Unruh~\cite{unruh}). At a horizon of surface
gravity $\kappa$ it is $4\pi^2$ times the Hawking temperature $k_BT_H=\hbar\kappa/2\pi c$; for a Schwarzschild hole $\kappa=c^4/4GM$.
\end{proposition}

\begin{proof}
With dimensions $[g]=LT^{-2}$, $[h]=ML^2T^{-1}$, $[c]=LT^{-1}$, the monomial $g^ah^bc^k$ has the dimension of energy, $ML^2T^{-2}$, only
for $(a,b,k)=(1,1,-1)$ \chk{Q1}. With $h=2\pi\hbar$, $E_{\text{ex}}/k_BT_U=4\pi^2$.
\end{proof}

The uniqueness means the form of~\eqref{eq:Eex} is not itself evidence; quantum field theory supplies the factor and the meaning. At standard
gravity $T_U=4.0\times10^{-20}$\,K; at the horizon of a solar-mass hole $T_H=6.2\times10^{-8}$\,K \chk{Q1}. \emph{Verdict: holds as a
reading} of the Unruh and Hawking temperatures.

\subsection{Time: the lapse}

In the GUID reading a clock asserts its existence at a rate that falls where the existence field is strong. In general relativity that
rate is the lapse, $N=\sqrt{1+2\Phi/c^2}$ in a weak field and $N=\sqrt{1-2GM/rc^2}$ outside a mass, which vanishes at the horizon
\chk{Q2}. It reproduces the redshift over the 22.5\,m Pound--Rebka tower, $2.46\times10^{-15}$~\cite{pr60}, and the rate difference of two
optical clocks 33\,cm apart, $3.6\times10^{-17}$~\cite{chou}. \emph{Verdict: holds}; it adds nothing to general relativity.

\subsection{Einstein's equation as an equation of state}

Jacobson derived Einstein's equation from thermodynamics~\cite{jacobson95}, and GUID's two thermal inputs are exactly his. Take any point
$p$ and any local Rindler horizon through it, generated by a null vector field $k^a$ with affine parameter $\lambda$, vanishing expansion and
shear at $p$, and approximate boost Killing field $\chi^a=-\kappa\lambda k^a$. Assume:
\begin{enumerate}[nosep]
\item the horizon has the Unruh temperature $T=\hbar\kappa/2\pi$ (units $c=k_B=1$), Proposition~\ref{pr:unruh};
\item its entropy is proportional to its area, $S=\eta A$, the boundary's floor (Hypothesis~\ref{h:horizon});
\item the Clausius relation $\delta Q=T\,\dd S$ holds for the energy flux across it.
\end{enumerate}
The heat flux is $\delta Q=\int T_{ab}\chi^a\dd\Sigma^b=-\kappa\int\lambda\,T_{ab}k^ak^b\,\dd\lambda\,\dd\mathcal A$. The area change follows
from the Raychaudhuri equation, $\dd\theta/\dd\lambda=-R_{ab}k^ak^b$ to first order, as $\delta A=-\int\lambda R_{ab}k^ak^b\,\dd\lambda\,\dd\mathcal A$.
Equating $\delta Q=(\hbar\kappa/2\pi)\eta\,\delta A$ for every $p$ and every null $k^a$ gives
$R_{ab}+f g_{ab}=(2\pi/\hbar\eta)T_{ab}$ for some function $f$; the contracted Bianchi identity and $\nabla^aT_{ab}=0$ fix
$f=-R/2+\Lambda$, so
\begin{equation}\label{eq:jacobson}
G_{ab}+\Lambda g_{ab}=\frac{2\pi}{\hbar\eta}\,T_{ab},\qquad G=\frac{1}{4\hbar\eta}\ \ (c=1).
\end{equation}
Newton's constant is set by the entropy per area, $\eta=1/4L_P^2$ \chk{Q1}, and the cosmological constant enters as an integration
constant. Jacobson's own conclusion fits GUID's classical frame: on this view, ``it may be no more appropriate to quantize the Einstein
equation than it would be to quantize the wave equation for sound in air''~\cite{jacobson95}. Two limits: DMP supplies neither the value of $\eta$ nor
a manifold, so the derivation is borrowed, not produced; and the entanglement-equilibrium version~\cite{jacobson16} needs a state, not a
schedule, which the axioms do not define. \emph{Verdict: holds as a reading}; the constant is open.

\subsection{Information radiation and the floor}

DMP's T31 says that a containing boundary emits only thermal radiation, carrying ``entropy and no information'', and that an outside reader
learns of the inside only its size and its staleness. With Hypothesis~\ref{h:horizon} that is the semiclassical horizon: a black hole of mass
$M$ radiates at $T_H$ and an outside observer knows its mass, spin and charge \tagP. Erasing a bit costs at least $k_BT\ln2$ of heat
(Landauer~\cite{landauer}); in the GUID reading entropy is the inefficiency of converting between existence and information, which is
Landauer's principle in other words. \emph{Verdict: holds as a reading.}

\section{Black holes}\label{sec:bh}

\subsection{The horizon as containment}

Hypothesis~\ref{h:horizon} maps DMP's containment onto the event horizon. The match is close. No act carries information across a containing
boundary (T31); no causal curve leaves the horizon \tagP. The counts every outside reader sees (Axiom~3) are the hair a stationary black hole
keeps: mass $M$, angular momentum $J$, charge $Q$ (Israel~\cite{israel}; Kerr~\cite{kerr}). Inside a Schwarzschild horizon the radial
coordinate is timelike and every future-directed path ends at the singularity: everything inside derives its causal future from it, a
local ground in the sense of G0. Whether crossing is survivable is answered by the equivalence principle: at the horizon of a
supermassive hole the tidal stretch across a body is small; for a stellar-mass hole it is lethal before the horizon.

\subsection{The containment split gives Christodoulou's mass formula}

DMP splits a contained boundary into what is inside and the counts every outside reader sees \tagD. Give the contained part a mass
$M_{\text{irr}}$ that sets a Schwarzschild boundary of radius $2GM_{\text{irr}}/c^2$; let the spin count $J$ be carried at that boundary as
a momentum $J/(2GM_{\text{irr}}/c)$ (an identification, \tagH); and combine the two as energy and momentum combine in special relativity. In
units $G=c=1$,
\begin{equation}\label{eq:christo}
M^2=M_{\text{irr}}^2+\left(\frac{J}{2M_{\text{irr}}}\right)^2 .
\end{equation}
This is Christodoulou's mass formula for the Kerr black hole~\cite{christodoulou}, with $M_{\text{irr}}=\sqrt{A/16\pi}$ \tagP; solved for
$M_{\text{irr}}$ it returns Kerr's horizon area for every spin $a=J/M^2\in[0,1]$, has no real solution above $a=1$ (the spin bound), and
gives the extractable share $1-M_{\text{irr}}/M=1-1/\sqrt2\approx29\%$ at maximal spin \chk{Q10}. The derivation uses DMP's split, two
physical inputs and one identification, and gives no number general relativity does not. \emph{Verdict: holds as a rederivation.}

\subsection{What a black hole gives up: Hawking--Penrose--Information radiation}

The first law of black-hole mechanics (Bardeen, Carter and Hawking~\cite{bch}) reads, with $G=c=1$,
\begin{equation}\label{eq:firstlaw}
\dd M=\frac{\kappa}{8\pi}\,\dd A+\Omega_H\,\dd J+\Phi_H\,\dd Q ,
\end{equation}
with $\Omega_H$ the angular velocity and $\Phi_H$ the electric potential of the horizon. It separates the two ways a black hole loses mass.
The rotational and electric terms can be extracted from outside at fixed or growing area: Penrose's process with particles~\cite{penrose69},
superradiant scattering with waves, and the Blandford--Znajek process with magnetic fields~\cite{bz}, the leading model of the jets of
active galaxies. All of them work in the ergosphere, which reaches $r=2GM/c^2$ at the equator and lies outside the horizon
$r_+=(GM/c^2)(1+\sqrt{1-a^2})$ for every spin \chk{Q10}. The area term is released only quantum mechanically, as Hawking's thermal
radiation, which lowers $A$.

\begin{definition}[Hawking--Penrose--Information radiation]
The \emph{Hawking--Penrose--Information (HPI) radiation} of a black hole is everything it gives up: the extraction of its rotational and
electric energy from the ergoregion (the $\Omega_H\dd J$ and $\Phi_H\dd Q$ terms of~\eqref{eq:firstlaw}) and its thermal Hawking radiation (the
$\kappa\,\dd A/8\pi$ term).
\end{definition}

In GUID terms the extraction is the existence frame being dragged along by the spinning hole, outside the horizon, so no act crosses it; the
Hawking part is the floor of T31. What leaves carries the counts and no content, as Axiom~3 requires. Spinning holes are more efficient
engines: the binding energy radiated by matter falling to the innermost stable orbit is $1-\sqrt{8/9}=5.7\%$ of its rest energy without
spin, $32\%$ at Thorne's limit $a=0.998$, and approaches $1-1/\sqrt3=42.3\%$ at maximal spin (Bardeen, Press and Teukolsky~\cite{bpt};
\chk{Q10}). At fixed mass every additional count (spin, charge) shrinks the area, so the more an outside reader knows of a boundary, the
less entropy it holds and the more of its energy can be extracted \tagP. The Hawking part is unobservable for astrophysical holes: a
solar-mass hole is at $6.2\times10^{-8}$\,K, colder than the $2.7$\,K microwave sky, so it absorbs more than it emits; only holes lighter
than about $0.6$ lunar masses are hotter than the sky, and a solar-mass hole in empty space would take more than $10^{66}$ years to
evaporate \chk{Q11}. \emph{Verdict: holds}; the Hawking part is \emph{untested in nature}.

\subsection{The information problem}

If evaporation is unitary, the entropy of the radiation first rises and then falls after the Page time~\cite{page93}; the island formula,
derived from the gravitational path integral with replica wormholes, reproduces that curve~\cite{penington,aemm}. Those derivations sum
over geometries, which is quantum gravity in the sense GUID's frame excludes. On the other side, a classical-quantum dynamics decoheres the
quantum system (Corollary~\ref{co:diffuse}), so it is not unitary on the quantum system alone and does not by itself produce a Page curve.
DMP's T31 as written, ``what still leaves is entropy and nothing else'', and its T27, which destroys contained items at their scheduled date,
side with the thermal picture. In GUID the question is therefore open in a specific form: whether the content held inside a horizon (which
may be quantum, Theorem~\ref{th:host}(ii)) returns in the radiation as content, while the frame stays classical. No observation reaches it.
\emph{Verdict: open}, in physics and in GUID.

\section{Holography and quantum error correction}\label{sec:holo}

\subsection{The strongest assembly}

The axioms can be assembled as a holographic quantum error-correcting code on a fixed background, the kinematics of the HaPPY
model~\cite{happy}. A part (``no further statements about its constituents'') becomes one block of the five-qubit code~\cite{laflamme},
which encodes one logical qubit in five physical ones with distance 3. Its encoding map is a \emph{perfect tensor} with six legs: every
bipartition of the six legs into three and three is a unitary map, so every three-leg cut carries three qubits of maximal entanglement
\chk{A1}. A layer, which Axiom~3 makes a tree rooted at its Duality, becomes a rooted tree of such tiles, one root and five children with
twenty boundary legs, each tile's logical leg the information sitting there; both tile maps are isometries \chk{A2}.

\paragraph{Minimal cuts.} For a boundary region $R$, the entropy of the reduced boundary state equals the number of legs cut by the minimal
cut through the tree that separates $R$ from its complement \chk{A2}. This is the Ryu--Takayanagi formula~\cite{rt},
$S(R)=|\gamma_R|/4G\hbar$ with $\gamma_R$ the minimal bulk surface anchored on $\partial R$, on a tree. DMP's ``fewest hops in total''
selects the same minimal surfaces.

\paragraph{Reconstruction.} The five-qubit code corrects any two erasures, so the logical qubit is recoverable from any three of the five
physical qubits and from no two: a $((3,5))$ quantum threshold scheme~\cite{cgl}. In mutual information with a reference $R$ purifying the
root, $I(R:\text{pair})=0$ for all ten pairs and $I(R:\text{triple})=2$ bits for all ten triples \chk{A3}. This is entanglement-wedge
reconstruction~\cite{adh}: a bulk operator is reconstructable on any boundary region whose wedge contains it.

\paragraph{The two witnesses.} The trust half and the information half each measure the stabiliser syndrome of the code. Stabilisers commute
with the logical operators, so the syndrome says that content is present and in the code, never what it is: this is Axiom~3's existence
clause, exactly. Both halves read the same syndrome; the gate passes only on two trivial syndromes; each of the 15 single-qubit faults gives
a distinct syndrome on which both refuse, and correction restores the content with fidelity 1. If the two halves instead read the content
itself, the content is disturbed (mean fidelity $0.50$), so ``full reach over the act'' must mean full reach over its existence \chk{A4}.

\paragraph{Two layers in one object.} Two parts whose content is a logical Bell pair reach the Tsirelson value $2\sqrt2\approx2.83$ of the
CHSH combination on their content, while the gates' verdicts, deterministic syndrome readings, give exactly the classical value 2
\chk{A5}. The gate is classical and the content quantum, at the same time and in the same assembly.

\subsection{Where the assembly stops: the classical centre}

In any code with complementary recovery, Harlow proved the quantum Ryu--Takayanagi formula
\begin{equation}\label{eq:harlow}
S(\rho_R)=\Tr(\rho\,\mathcal L_R)+S_{\text{bulk}}(\rho,\mathcal M_R),
\end{equation}
where $\mathcal M_R$ is the algebra of bulk operators reconstructable on $R$ and the \emph{area operator} $\mathcal L_R$ lies in the centre of
$\mathcal M_R$: it commutes with everything either side can do~\cite{harlow}. DMP's structures recomputed by both halves are such central
elements. Lemma~\ref{lem:classical} then says that in GUID the area is a number, never an operator with a spread of values: the axioms
allow one fixed-area sector at a time. In holography fixed-area states are a special, though useful, class (Dong, Harlow and
Marolf~\cite{dhm}); the general state of quantum gravity superposes areas. So the assembly reproduces the kinematics of holographic codes on
a fixed background (\chk{A1}--\chk{A5}) and stops exactly at superpositions of geometry.

\paragraph{A dissenting reading.} One of the review passes behind this work scored the same facts the other way. Split the geometry, as the
background-field method does: the background is a DMP structure and classical; the linearised fluctuation, which carries the Newtonian
pull between laboratory masses, is content between parts, and linearised quantum gravity on a fixed background already predicts the
entanglement of the masses~\cite{christodoulou23}. On that reading DMP's classical structures are the part of quantum gravity that is
classical in quantum gravity too, the centre of~\eqref{eq:harlow}, and the linearised field is quantum content on that classical
centre. The two readings agree on every fact: the shared structures are classical, quantum content including a linearised field can
move between parts by grant and stay entangled while every act is recorded, the axioms derive no constant, and they forbid coherent
superpositions of backgrounds while their records exist. They differ in one place: whether the linearised field that pulls laboratory
masses belongs to the frame (the proposed theory, which predicts $\mathcal N=0$ in Section~\ref{sec:exp}) or to the content (the dissent,
which predicts $\mathcal N>0$). DMP hosts either; the experiment decides between them.

\paragraph{The tree is not our space.} Read as space, a tree's boundary is totally disconnected (the $p$-adic line of $p$-adic
AdS/CFT~\cite{hmss,gubser}), and no tree approximates a plane: an $n$-cycle needs distortion at least $n/3-1$ in any tree~\cite{rr}. So a
literal tree cannot be our space. Keeping the hierarchy as the routing tree of each request over a graph with loops removes this objection
and touches none of the results above. Read as the tree of histories, each branch one consistent record (decoherent
histories~\cite{gellmannhartle}), the tree owes no spacelike structure at all.

\section{Lower dimensions}\label{sec:21}

Why 3+1? In $D$ spacetime dimensions a massless spin-2 field has $D(D-3)/2$ polarisations and the Weyl tensor, the part of the curvature
not fixed by local matter, has $D^2(D^2-1)/12-D(D+1)/2$ components: two and ten in 3+1, none and none in 2+1 \chk{Q8}. In 2+1 a mass does not
attract another at all; it cuts a wedge out of space (Deser, Jackiw and 't~Hooft~\cite{djt}). Gravity there is, perturbatively, a
Chern--Simons gauge theory~\cite{witten88}, with a holographic boundary algebra for negative cosmological constant~\cite{bh}; whether pure
quantum gravity in 2+1 exists as a consistent theory is open~\cite{mw}. Its degrees of freedom are global: the holonomies of the flat
connection around particles and handles, which do not commute once quantised~\cite{nr}, and point particles scatter through their conical
deficits~\cite{thooft88}.

So 2+1 is not an exception to Theorem~\ref{th:host}. With no local field there is no local mediator to forbid, but the global holonomies are
quantum, and a DMP structure that carries them is broadcast and so can carry one classical sector of them, never their superposition; with
matter, the same rule decoheres the particles' relative positions. In GUID's terms, the existence frame in 2+1 would have nothing local
that could be held as a classical record, and what it does carry is global and quantum. DMP can run 2+1 quantum gravity as content, but
its frame cannot be 2+1 gravity;
and our universe is 3+1, whose local degrees of freedom, gravitational waves, are measured~\cite{gw150914}. \emph{Verdict: derived}, with
the dimension count as \emph{holds}.

\section{Cosmology}\label{sec:cosmo}

\subsection{The beginning: the start}

G0 is stated, not derived, and the minting is the first act: before it there is no act and therefore no record \tagD.
Lemma~\ref{lem:classical} rests on records, so it applies only once acts exist. GUID describes the beginning as one event, the
\emph{start}. For anything to begin to exist it must be possible for it to exist, and the first step of existing is to inflate, through the
decay of the oscillons that coalesce once the initial energy has been emitted. Inflation and that decay cannot be separated or ordered,
because there is no causality yet in which to order them; together they are the start.

\begin{hypothesis}[the start]\label{h:minting}
Inflation and the decay of the oscillons it ends in are one event, the start, which is the minting. It has no internal order the axioms can
refer to, and the classical era of Lemma~\ref{lem:classical}, with causality in DMP's sense (acts between existing parts), begins as it
completes. \tagH
\end{hypothesis}

The physics of the start is standard \tagP. Inflation ends as the inflaton oscillates about the minimum of its potential; in a broad class
of models self-resonance fragments it into oscillons, long-lived localised lumps of the field, which dominate the universe for a while
(Amin et al.~\cite{amin}); their decay into the particles of the Standard Model reheats the universe (Kofman, Linde and
Starobinsky~\cite{kls}). In physics these are stages of one field's history; GUID's statement that they form one event agrees with that,
since the axioms, which need acts to order anything, have nothing by which to order them. Within the start the axioms place no restriction
on the geometry, which may be in superposition; after it, every geometry the axioms define is recorded and classical. This matches the
standard semiclassical picture, in which the geometry of the universe becomes classical through decoherence at its very beginning
(Halliwell~\cite{halliwell}; Kiefer~\cite{kiefer87}), and the no-boundary proposal, in which there is no ``before'' (Hartle and
Hawking~\cite{hartlehawking}). In GUID the start is therefore the boundary of quantum gravity in time: whatever is quantum about the
geometry itself lies within it, where DMP by its own words makes no statement.

\begin{vnew}
\begin{hypothesisprime6}[the strong split; refines Hypothesis~\ref{h:minting}]
\vtag{} The start also contains the split of the strong force from the electroweak force, and causality begins with that split. If the two
were ever one force, the unification lies before time: within causality the two were never one. \tagH
\end{hypothesisprime6}

\noindent\vtag{} The Standard Model agrees within causality: in it the strong and the electroweak force are separate gauge interactions
and never meet; the three couplings come close near $10^{14}$ to $10^{16}$\,GeV but do not meet at one point~\cite{amaldi}. A unification
would leave three traces in the causal era, and they make the hypothesis testable in part \tagP. (1) \emph{Magnetic monopoles.} The split of
a unified force generically creates very heavy magnetic monopoles, about one per causal volume (Preskill~\cite{preskill}); had the split
happened after inflation, the universe would be full of them. None has been seen: MACRO bounds their flux below
$1.4\times10^{-16}$\,cm$^{-2}$\,s$^{-1}$\,sr$^{-1}$~\cite{macro}. The split must therefore happen before or during inflation, and reheating
must stay below the unification energy, which was one of Guth's reasons for inflation~\cite{guth}. Hypothesis~6$'$ puts it exactly there:
the ordering is not only allowed, it is required. (2) \emph{Proton decay.} The heavy gauge bosons of a unified force, of mass $M_X$, still
let protons decay after the split, with a lifetime of order $\tau\sim M_X^4/(\alpha_G^2m_p^5)$ to within a factor of about ten: about
$5\times10^{35}$ years for $M_X=10^{16}$\,GeV and $\alpha_G=1/40$ \chk{Q20}. Super-Kamiokande finds $\tau(p\to e^+\pi^0)>2.4\times10^{34}$
years~\cite{superk}, which requires $M_X\gtrsim5\times10^{15}$\,GeV in this estimate and excludes the simplest unified theory of Georgi and
Glashow~\cite{georgiglashow}; Hyper-Kamiokande will reach much further~\cite{hyperk}. (3) \emph{The couplings.} A real unification needs
particles beyond the Standard Model that make the couplings meet. The bound on inflation, $V^{1/4}<1.4\times10^{16}$\,GeV
(Section~\ref{sec:tensor}), allows a split at the unification scale within the start \chk{Q20}. \emph{Verdict: holds as a reading}; the
ordering is required by the absence of monopoles; the unification itself is \emph{untested}, with proton decay as its live test.

\paragraph{\vtag{} The timeline.} Read together, Hypotheses~6, 6$'$ and Section~\ref{sec:higgsplit} give one sequence:
\begin{enumerate}[nosep]
\item \emph{The start.} The strong force splits from the electroweak force; inflation dilutes the monopoles of the split; the decay of the
oscillons reheats the universe to below the unification energy. Causality begins.
\item \emph{About $10^{-11}$\,s, at 160\,GeV.} The symmetric electroweak state becomes unstable; the Higgs vacuum takes over granting rest
mass to the W, the Z and the matter particles, and the information field, the photon and with it $\alpha\approx1/128$, splits out of the
electroweak field.
\item \emph{Afterwards.} $\alpha$ runs down to $1/137$ at low energies; nucleosynthesis, recombination and the late universe follow.
\end{enumerate}
\end{vnew}


\subsection{The universe at its own Schwarzschild threshold}

For a flat Friedmann universe with Hubble rate $H$, the mass inside the Hubble radius, $M_H=\tfrac43\pi(c/H)^3\rho_c$ with
$\rho_c=3H^2/8\pi G$, has a Schwarzschild radius equal to the Hubble radius,
\begin{equation}\label{eq:kugel}
\frac{2GM_H}{c^2}=\frac cH ,
\end{equation}
exactly and at every epoch \chk{Q12}. This is the coincidence behind Pathria's ``universe as a black hole''~\cite{pathria}. GUID reads the
early universe as a Kugelblitz, a horizon formed by energy alone, that could not contain its own existence; the reading holds in the sense
of~\eqref{eq:kugel} and at every epoch, not only the earliest, and the universe is not a black hole in the sense of general relativity,
having no outside. \emph{Verdict: holds as a reading.}

\subsection{Inflation as the possibility space of existence}

In GUID, inflation is the space that existence makes for itself wherever existence is possible. Physically it is a period of nearly
exponential expansion, $a\propto e^{Ht}$, which dilutes a spatial curvature $|\Omega-1|$ by $e^{-2N}$ over $N$ e-folds (Guth~\cite{guth})
\tagP. Afterwards the curvature share grows again, as $a^2$ while radiation dominates and as $a$ while matter does; from a reheating
temperature of $10^{15}$\,GeV to today it grows by about $5\times10^{51}$, so an order-one initial curvature needs about 62 e-folds to meet
Planck's bound $|\Omega_k|\lesssim5\times10^{-3}$~\cite{planck18} \chk{Q13}. \emph{Verdict: holds as a reading}; it supplies the motive, not
the number.

\subsection{Primordial gravitational waves}\label{sec:tensor}

During inflation the vacuum fluctuations of the metric are stretched into a nearly scale-invariant background of gravitational waves
(Starobinsky~\cite{starobinsky79}). Their power spectrum is
\begin{equation}\label{eq:tensor}
\mathcal P_t=\frac{2}{\pi^2}\left(\frac{\hbar H}{M_{\text{pl}}c^2}\right)^2=\frac{16\,G\hbar H^2}{\pi c^5},
\end{equation}
with $M_{\text{pl}}^2=\hbar c/8\pi G$, linear in $\hbar$ and zero for a classical metric \chk{Q14}; its ratio to the scalar spectrum is
$r=16\epsilon$ in single-field slow roll, with the consistency relation $n_t=-r/8$ for its tilt. Because~\eqref{eq:tensor} vanishes as
$\hbar\to0$, Krauss and Wilczek argued that detecting this background would establish that the gravitational field was quantised during
inflation~\cite{krausswilczek} \tagP. The current bound is $r<0.036$ at 95\% confidence (BICEP/Keck~\cite{bk18}; $r<0.032$ combined with
Planck and baryon acoustic oscillations~\cite{tristram}), which puts the energy scale of inflation below $V^{1/4}=1.4\times10^{16}$\,GeV
and its Hubble rate below $\hbar H=4.7\times10^{13}$\,GeV \chk{Q14}. Starobinsky's $R^2$ inflation~\cite{starobinsky80}, which fits the
scalar data, predicts $r\approx12/N^2\approx3\times10^{-3}$ at $N=60$ \chk{Q14}, the target of the next generation of polarisation
surveys~\cite{litebird}.

Under Hypothesis~\ref{h:minting} the tensor modes are generated within the start, where the axioms do not constrain the geometry,
and become classical, as the standard account says they do through squeezing and decoherence~\cite{polarski}, as the start completes. GUID
therefore allows the standard spectrum and its consistency relation, and a detection would be the first observation of the quantum
geometry of the start. GUID's own predictions about the geometry begin after the start. \emph{Verdict: live}; consistent with every bound.

\subsection{The late universe: the balanced state}

In the flat $\Lambda$CDM model that fits the data ($\Omega_m=0.315$, $\Omega_\Lambda=0.685$~\cite{planck18}), the Hubble rate falls at every
moment from the past into the far future and tends to the de Sitter rate $H_\infty=\sqrt{\Omega_\Lambda}\,H_0\approx0.83\,H_0$; the expansion
has accelerated since $z=(2\Omega_\Lambda/\Omega_m)^{1/3}-1\approx0.63$ \chk{Q13}. A universe settling into de Sitter expansion cools towards
an empty final state. GUID reads this as the cosmological balanced state, the equilibrium in which the need for further existence has
fallen to a constant. \emph{Verdict: holds as a reading} in $\Lambda$CDM. The newest data leave the end open: DESI's second data release
prefers a dark energy that weakens with time at 2.8 to 4.2 standard deviations, depending on the supernova sample~\cite{desi2}; if that holds,
the future is not exactly de Sitter and the reading must be redone.

\subsection{Dark matter: remnants trapped in the gravitational field}\label{sec:darkmatter}

\begin{hypothesis}[trapped remnants]\label{h:remnant}
The modelling of this paper explicitly allows remnants of interactions to be trapped in the gravitational field: energy and momentum that
interactions have left in the existence frame, or behind a boundary the frame has closed, continue to gravitate after the interaction.
\tagH
\end{hypothesis}

DMP motivates the hypothesis twice: a record outlives its act (Axiom~3), and a containing boundary keeps what it holds while every outside
reader sees only its counts (T31). Dark matter, as far as the data show, is matter of exactly that kind: it gravitates and does nothing
else that has been seen. Within the classical-frame modelling there are three routes by which remnants are trapped, and each can be
calculated.

\paragraph{The stochastic remnant.} Corollary~\ref{co:diffuse} makes the frame diffuse: every interaction of quantum matter with the
classical frame leaves a random kick in the field. Oppenheim and Russo computed the expectation value of the metric in classical-quantum
gravity from its path-integral formulation, for static, spherically symmetric configurations~\cite{oppenheimrusso}. Around a mass $M$ the
Newtonian potential acquires a linear and a quadratic term,
\begin{equation}\label{eq:stochpot}
\Phi(r)=-\frac{GM}{r}-\frac{\gamma_0}{2}-\frac{\gamma_1}{2}\,r-\frac{\gamma_2}{2}\,r^2 ,
\end{equation}
where the quadratic term is the cosmological constant of Schwarzschild--de Sitter space, $\gamma_2=\Lambda c^2/3>0$, which repels, and
the linear term is anticorrelated with it: its mean is negative, $\gamma_1<0$, so it attracts, with $|\gamma_1|$ of order $c^2\sqrt\Lambda$.
Physically, where accelerations are low the variance of the field's acceleration exceeds the square of its mean and acts as an additional,
entropic attraction. The circular speed is
\begin{equation}\label{eq:vcirc}
v^2(r)=r\,\frac{\dd\Phi}{\dd r}=\frac{GM}{r}+\frac{|\gamma_1|}{2}\,r-\gamma_2\,r^2 .
\end{equation}
The first two terms balance at $r_*=\sqrt{2GM/|\gamma_1|}$, where
\begin{equation}\label{eq:btf}
v_*^4=2GM|\gamma_1| .
\end{equation}
Around $r_*$ the rotation curve is roughly flat, changing by about 12\% between $r_*/2$ and $2r_*$, and its height obeys $v^4\propto M$,
the form of the observed baryonic Tully--Fisher relation. For a galaxy of $10^{11}$ solar masses in baryons and $|\gamma_1|=10^{-10}$\,m/s$^2$,
the order Oppenheim and Russo find, $r_*=17$\,kpc and $v_*=227$\,km/s, with the $\Lambda$ term below $10^{-4}$ of $v_*^2$ there \chk{Q18}.
The acceleration scale is tied to the cosmological constant: $c^2\sqrt{\Lambda/3}/2\pi=8.6\times10^{-11}$\,m/s$^2$, close to Milgrom's
empirical $a_0=1.2\times10^{-10}$\,m/s$^2$~\cite{milgrom} \chk{Q18}. On this route a single mechanism, the frame's diffusion, gives both a
dark-energy term and a dark-matter-like term. The calculation is preliminary, as its authors stress: it is static and spherically
symmetric, depends on an infrared regularisation, fixes $\gamma_1$ by boundary conditions rather than by a coupling constant, and does not
yet address gravitational lensing or the microwave background. Fixing the boundary condition by GUID's outermost horizon gives a value
eight times too strong in $v^4$ (Section~\ref{sec:bridges}). A deterministic relative of the same idea exists in tetrad gravity: Hehl
and Mashhoon found that a memory of the frame, a nonlocal dependence on its past, acts in the linear approximation as dark matter given by
an integral transform of the matter~\cite{hehlmashhoon}.

\paragraph{The contained remnant.} The second route is containment (Hypothesis~\ref{h:horizon}). Energy that collapsed behind a horizon
during the start gravitates through its counts alone and is dark by T31: primordial black holes. The start itself makes them. When the
inflaton fragments into oscillons, which in a broad class of models dominate the universe for a while after inflation~\cite{amin}, enough
of them can collapse into black holes of sublunar mass to make up all of the dark matter (Cotner, Kusenko and Takhistov~\cite{cotner}).
Primordial black holes can be all of the dark matter only in the asteroid-mass window, about $10^{14}$ to $10^{19}$\,kg ($10^{17}$ to
$10^{22}$\,g), which current constraints leave open~\cite{gortongreen}. The window's light edge is set by the Hawking part of their HPI
radiation: a hole of $10^{14}$\,kg glows at $1.2\times10^9$\,K, $k_BT\approx0.1$\,MeV, and its evaporation products (MeV gamma rays, and
positrons that annihilate at 511\,keV) bound how many such holes there can be; holes lighter than a few $10^{11}$\,kg have evaporated
already \chk{Q18}. At the heavy edge, $10^{19}$\,kg with a Schwarzschild radius of 15\,nm and a temperature of $1.2\times10^4$\,K,
microlensing takes over \chk{Q18}. The prediction of this route is specific: if dark matter is the contained remnant of the start, the
Hawking part of HPI radiation, unobservable for astrophysical black holes (Section~\ref{sec:bh}), becomes observable at the window's light
edge.

\paragraph{The radiated remnant.} Gravitational waves are remnants of interactions that propagate in the frame. Pulsar timing has found a
nanohertz background with characteristic strain $h_c\approx2.4\times10^{-15}$ at the frequency $1/\text{yr}$ (NANOGrav~\cite{nanograv}),
which corresponds to $\Omega_{\text{gw}}=2\pi^2f^2h_c^2/3H_0^2\approx8\times10^{-9}$ per logarithmic frequency \chk{Q18}: more than $10^7$
times too little for the dark matter, and as radiation ($w=1/3$) it does not cluster. It is a real remnant trapped in the field, but not
the dark matter.

\paragraph{What every remnant must satisfy.} Four measurements constrain every route \tagP. (1) It must already be there at recombination:
the acoustic peaks of the microwave background measure the dark matter density at $z\approx1100$, before the first star, as
$\Omega_ch^2=0.120$, the same comoving value later structure requires, with the baryons then already at the share
$\Omega_b/\Omega_m=0.157$~\cite{planck18} \chk{Q15}. So the remnant formed during the start, not from later acts; and the baryon content
was fixed early, as nucleosynthesis and the microwave background agree. (2) It must be collisionless and move with the galaxies: in the
merging cluster 1E0657$-$558 the lensing mass follows the galaxies, not the hot gas that holds most of the baryons (Clowe et
al.~\cite{clowe}). (3) It cannot be the known neutrinos: they are less than $0.7\%$ of the dark matter density, since cosmology bounds
their mass sum below $0.072$\,eV~\cite{desi}; a dwarf galaxy made of them would need each to weigh more than about 170\,eV by the Pauli
principle (Tremaine and Gunn~\cite{tremainegunn}); and relic neutrinos move at about 3200\,km/s and stellar ones at nearly $c$, against a
Galactic escape speed of 533\,km/s~\cite{piffl} \chk{Q15}. (4) It must leave $\alpha$ unchanged to about a part per million over ten
billion years (Section~\ref{sec:alpha}).

\emph{Verdict: live.} The contained route passes all four requirements by construction; the stochastic route reproduces the shape and
scale of galactic rotation curves in a preliminary calculation and must still pass lensing, the acoustic peaks and the merging cluster;
the radiated route is real but far too small.

\subsection{The fine-structure constant as the information field}\label{sec:alpha}

In GUID $\alpha$ is set by the information scalar $\phi$ (Hypothesis~\ref{h:info}), so it may vary where $\phi$ does. The data are
tight \tagP. Optical clocks bound today's drift to $\dot\alpha/\alpha=(1.0\pm1.1)\times10^{-18}$ per year~\cite{lange}; quasar absorption
spectra bound changes over ten billion years to about a part per million~\cite{espresso}; the microwave background bounds the value at
recombination to a few parts per thousand~\cite{planckxxiv}. MICROSCOPE bounds the scalar's coupling to matter by~\eqref{eq:lambdabound}.
GUID's forward statement is that $\alpha$ does not grow: its drift is zero or negative. A measured positive drift would refute it.
\emph{Verdict: live.} 

\begin{vnew}
\vtag{} GUID's earlier statement that $\alpha$ began at one is withdrawn as a physical claim. The forced dissipation of G0 is no act and
precedes any measurement; GUID labels it one axiomatically, and the label has no value in calculations. Physically $\alpha$ does not exist
before the electroweak split: there is no photon then, only the couplings of the electroweak field, and $\alpha$ is born at the split, near
$1/128$ at the Z mass, from which it runs down to $1/137$ at low energies (Section~\ref{sec:higgsplit}; \chk{Q16}, \chk{Q20}). Something
close to the present value came out as the fields separated, which is what happened. \emph{Verdict: holds.} (Version 5 had tested the
physical form of the statement and found it beyond known physics: the Standard Model gives $1/\alpha\approx105$ even at the Planck energy
\chk{Q16}.)
\end{vnew}

\section{Matter in the existence frame}\label{sec:matter}

\subsection{The spin ladder}

GUID reads spin as the rungs of the Duality, and the Standard Model fills each rung \tagP:

\begin{center}\small
\begin{tabularx}{\linewidth}{@{}lLLL@{}}
\toprule
Spin & GUID reading & Particles & Physics \\
\midrule
0 & existence itself & Higgs field; the scalar $\phi$ & the Higgs vacuum value gives rest mass; $\phi$ sets $\alpha$ \\
1/2 & a duality, which can stay in balance & quarks, leptons & a massive Dirac fermion is two chiral halves joined by its mass \\
1 & information itself; delivered, an act & photon, gluons, W, Z & the carriers of every interaction \\
2 & the perfect duality, two in one & the frame's wave & a massless spin-2 field must couple to all energy: general relativity~\cite{weinberg64,deser70} \\
\bottomrule
\end{tabularx}
\end{center}

\noindent The rungs are told apart by rotation: a state of helicity $s$ picks up the phase $e^{is\theta}$ under a turn by $\theta$ about its
momentum, so a spin-1 quantum looks the same again after $360^\circ$, a spin-2 quantum after $180^\circ$, and a spin-1/2 state only after
$720^\circ$, returning with a minus sign after $360^\circ$. \emph{Verdict: holds as a reading.}

\subsection{The Higgs vacuum as the grant of existence}

In GUID existence is granted by the ground. The Higgs field is the one field whose value in empty space is not zero: with potential
$V(H)=-\mu^2H^\dagger H+\lambda_H(H^\dagger H)^2$ it rests at $\langle H^\dagger H\rangle=v^2/2$, $v=\mu/\sqrt{\lambda_H}=246$\,GeV, and an
elementary particle has its rest mass by its coupling to that background~\cite{higgs64,be64} \tagP. GUID reads the background value, which
quantum field theory treats as a classical number, as G0's grant of rest mass, on the existence side and classical, in agreement with
Lemma~\ref{lem:classical}; the Higgs boson, a quantum ripple of the field that decays in about $10^{-22}$\,s, is content. The vacuum at the
minimum of its potential is the physical counterpart of the balanced state $\Phi=0$ of~\eqref{eq:vibe}. Like G0 (``stated, not
derived''), the Higgs potential is an input of the Standard Model, not a derivation.

The reading has one sharp limit, and it is gravitational. The Higgs field gives only part of the mass of matter: about 9\% of a nucleon's
mass is the quark mass terms; the rest is the energy of quark motion and of the gluon field~\cite{yang}. Gravity weighs all of it: a gravity
that weighed only the Higgs-given share would make titanium and platinum fall differently by $10^{-3}$, against MICROSCOPE's
$10^{-15}$ \chk{Q6}. So the existence frame responds to all existence, granted (the Higgs share) and bound (the internal energy of
hadrons), and the Higgs field is not the gravitational field. And gravity is far too weak to give mass: its coupling to the electron,
$Gm_e^2/\hbar c=1.75\times10^{-45}$, and its share of the binding of hydrogen, $4.4\times10^{-40}$ of the electric one, are negligible
\chk{Q16}. \emph{Verdict: holds as a reading} of the Higgs vacuum as the grant; \emph{refuted} as an identity of the Higgs field with
gravity.

\begin{vnew}
\subsubsection*{\vtag{} The Higgs vacuum as the existence that separated out}\label{sec:higgsplit}

In version 5.1 GUID makes the Higgs field the leading candidate for the existence field, the one that separated out when the fields did.
Physics says when \tagP. Above about 160\,GeV the Higgs field rests at zero, and no elementary particle has rest mass; lattice calculations
find that, with the measured Higgs mass, the change is a smooth crossover at $T_c=159.5\pm1.5$\,GeV (D'Onofrio and
Rummukainen~\cite{donofrio}), about $10^{-11}$\,s after the start \chk{Q20}. GUID's posit, in its form that holds, has four parts.

\emph{Before the split, existence asserted itself through interaction with existence.} In the hot plasma before the crossover the particles
carry only thermal masses, of order coupling times temperature, which come entirely from their interactions with everything around them.
\emph{Verdict: holds as a reading.}

\emph{That became unstable.} As the universe cooled, the curvature of the Higgs potential at zero field turned negative and the field
rolled to its present value $v=246$\,GeV; with a 125\,GeV Higgs this is gradual, not a sudden phase transition. \emph{Verdict: holds}, as a
crossover.

\emph{The information field split out of the electroweak field.} At the split the photon emerges as the fixed mixture
$A=\sin\theta_W\,W^3+\cos\theta_W\,B$ of the two electroweak gauge fields, and the electric charge is born as $e=g\sin\theta_W$, so that $\alpha$
appears near $1/128$ \chk{Q20}. \emph{Verdict: holds}, exactly.

\emph{The Higgs field became needed then.} At that moment the W and Z take their masses from the vacuum, $m_W=gv/2\approx80$\,GeV and
$m_Z=m_W/\cos\theta_W\approx91$\,GeV within half a per cent of the measured values at tree level \chk{Q20}, while the photon stays massless:
the Higgs vacuum grants existence to the bosons that need it at the very moment the information field emerges. \emph{Verdict: holds.} The
Standard Model Higgs belongs to this split, the weak from the electromagnetic; a Higgs-type field for the earlier split of the strong force,
if that split was a unification breaking (Hypothesis~6$'$), would be a second, hypothetical field.

\emph{An open tension.} In Higgs inflation the same field also drives the start: with a coupling $\xi H^\dagger HR$ of the Higgs to
curvature, $\xi$ of order $10^4$, it fits the observed spectrum (Bezrukov and Shaposhnikov~\cite{bezrukov}), which would make the
existence field the field of the start itself. That coupling is a scalar multiplying the curvature, which Proposition~\ref{pr:gamma}
excludes in this form; written in the Einstein frame the field is minimally coupled. Whether Axiom~1 forbids Higgs inflation or only one way
of writing it is open.
\end{vnew}

\subsection{The VIBE coefficient}

GUID names $\alpha$ the \emph{VIBE coefficient}, after the Variable Influence Balancing Equation~\eqref{eq:vibe}, and reads it as the strength
with which existence interacts with itself before it must re-ground itself. In quantum electrodynamics $\alpha$ is the strength with which
a charge interacts with the electromagnetic field and so, through it, with itself; each emission and reabsorption of a virtual photon costs
one power of $\alpha$. The cleanest measure is the electron's anomalous magnetic moment, whose first term $a_e=\alpha/2\pi$
(Schwinger~\cite{schwinger}) gives the measured value~\cite{fan} to $0.15\%$ \chk{Q16}. The same constant is the ratio of the electron's
three lengths, $r_e/\bar\lambda_C=\bar\lambda_C/a_0=\alpha$ \chk{Q16}. \emph{Verdict: holds as a reading}; the VIBE coefficient is a name for
$\alpha$, not a new constant.

\subsection{The forces as tiers of one interaction}

GUID reads the forces as tiers of one interaction: the strong force binds the cluster, the weak force keeps it stable, electromagnetism
links clusters of clusters. Each does that job \tagP: the strong force binds quarks into nucleons and nucleons into nuclei, the weak force
moves nuclei towards the valley of stability, electromagnetism binds atoms and molecules. The weak and electromagnetic forces are one
electroweak interaction~\cite{glashow,weinberg67,salam}. Whether all three meet at one energy is open: in the Standard Model the couplings
come close near $10^{14}$ to $10^{16}$\,GeV and do not meet at one point~\cite{amaldi}. Gravity is not among the tiers in GUID: it is the
frame the tiers act in. \emph{Verdict: holds as a reading}; unification \emph{untested}. \nw{\vtag{} In GUID's timeline the
strong force splits from the electroweak force at the start (Hypothesis~6$'$), and the electroweak field splits into the weak and the
electromagnetic force about $10^{-11}$\,s later, when the Higgs vacuum appears (Section~\ref{sec:higgsplit}).}

\section{Hosting: the Existence Turing machine}\label{sec:host}

Theorem~\ref{th:host} says that DMP, supplied with the base configuration of our universe, hosts the proposed theory, and its part~(ii)
that the information half can carry any quantum content. Assembled from DMP parts, a universal quantum computer
can run as content, and by the Church--Turing--Deutsch principle it approximates every finitely realisable physical system to any finite
precision~\cite{deutsch}. GUID calls such an assembly the \emph{Existence Turing machine}; a sealed file that carries the assembly's
recipe, its digest, the versions of the axioms and of the proof chain, and the verdict lets the same machine be redeployed and checked
anywhere. Two consequences matter for quantum gravity. First, every candidate theory of quantum gravity in Section~\ref{sec:landscape} can
also run on the machine with its predictions, so the axioms do not forbid quantum gravity in nature; a theory whose geometry is quantum runs
as content, not as DMP's frame. Second, the machine's own structure, the gates, records and layout, is classical, and in the proposed theory
it is the existence frame itself. A
check of the second statement, with three five-qubit parts whose record and counts are identical for all sixteen pairs of mass contents
while the contents end entangled, belongs to the dissenting review and rests on one construction choice (existence read per part, not per
mass position); we do not use it here.

\section{Relation to the approaches to quantum gravity}\label{sec:landscape}

\begin{center}\small
\begin{longtable}{@{}p{0.18\linewidth}p{0.22\linewidth}p{0.28\linewidth}p{0.22\linewidth}@{}}
\toprule
Approach & What is quantum & Relation to GUID & What tells them apart \\
\midrule
\endhead
Effective field theory~\cite{donoghue,bdh} & metric fluctuations below $E_P$ & GUID keeps the classical part; the graviton-loop term
of~\eqref{eq:eft} is absent & a term $10^{-70}$ of Newton's at 1\,m \chk{Q1}: unmeasurable \\
String theory, AdS/CFT~\cite{maldacena} & everything; geometry emerges from a quantum boundary & hostable as content; the code kinematics
of Section~\ref{sec:holo} are shared, with a classical centre & entanglement of masses \\
Loop quantum gravity~\cite{ashtekar,rovellismolin} & the triad/tetrad; area has a discrete spectrum & the frame is LQG's own variable, not
quantised; area a number & area spectrum (not accessible); BMV \\
Asymptotic safety~\cite{weinberg79,reuter} & the metric, with a UV fixed point & hostable & BMV \\
Causal dynamical triangulations~\cite{cdt} & sum over geometries & allowed within the start only & BMV \\
Causal sets~\cite{bombelli} & discrete order; $\Lambda\sim1/\sqrt N$ & DMP's order is a tree, not a sprinkled causal set & Lorentz
invariance needs random sprinkling \\
Thermodynamic gravity~\cite{jacobson95} & the microscopic entropy only & GUID's thermal reading, Section~\ref{sec:thermo} & none
beyond GR \\
Classical-quantum gravity~\cite{oppenheim,osswd} & matter only & GUID's dynamics, Corollary~\ref{co:diffuse} & diffusion against
decoherence \\
Channel models~\cite{ktm,tilloydiosi} & matter only & GUID's architecture, Proposition~\ref{pr:ktm} & noise and decoherence rates \\
Gravity-related collapse~\cite{diosi87,penrose96} & matter, collapsing & parameter-free version excluded~\cite{donadi}; GUID needs a free strength &
spontaneous radiation \\
Semiclassical gravity~\cite{moller,rosenfeld63} & matter; metric sourced by $\langle T\rangle$ & rejected: nonlinear; Page--Geilker &
done~\cite{pagegeilker} \\
Modified dynamics (MOND)~\cite{milgrom} & nothing; a law at low acceleration & the stochastic remnant gives a linear term of
size $a_0$, tied to $\Lambda$ (Section~\ref{sec:darkmatter}) & lensing, acoustic peaks, merging clusters \\
Spacetime from entanglement~\cite{vanraamsdonk,ermepr} & content; geometry emergent & GUID's frame as the classical centre of such a geometry is
open & unclear \\
\bottomrule
\end{longtable}
\end{center}

\section{Predictions}\label{sec:predictions}

\begin{center}\small
\begin{longtable}{@{}p{0.27\linewidth}p{0.22\linewidth}p{0.2\linewidth}p{0.23\linewidth}@{}}
\toprule
Observable & GUID & Quantised metric & Status \\
\midrule
\endhead
Entanglement of two masses through gravity alone~\cite{bose,mv} & none, $\mathcal N=0$ & $\mathcal N=\tfrac12|\sin(\Delta\varphi/2)|$ &
proposed; 25\,s at $10^{-14}$\,kg \chk{Q5} \\
LOCC inequality with massive oscillators~\cite{lpp} & satisfied & can be violated & proposed \\
Decoherence and diffusion~\cite{osswd} & both non-zero, product bounded by~\eqref{eq:tradeoff} & neither required & bounded from both sides \\
Universality of free fall & exact for the frame; a scalar term $\eta=2\lambda^2f_\oplus\Delta f$ & exact & $|\eta|\lesssim10^{-15}$~\cite{microscope} \\
Light bending $\gamma$ & exactly 1 & 1 & $\gamma-1=(2.1\pm2.3)\times10^{-5}$~\cite{cassini} \\
Speed of gravitational waves & $c$ & $c$ & $|c_{\text{gw}}/c-1|\lesssim10^{-15}$~\cite{gw170817} \\
Single-graviton detector clicks~\cite{tobar} & possible, from a classical wave & possible & do not discriminate~\cite{cdr} \\
Primordial tensor modes & standard spectrum, generated within the start & $r=16\epsilon$, $n_t=-r/8$ & $r<0.036$~\cite{bk18} \\
Drift of $\alpha$ & zero or negative & none predicted & $(1.0\pm1.1)\times10^{-18}$/yr~\cite{lange} \\
Hawking radiation & thermal to the outside reader; return of content open & Page curve & not observable \\
Galactic rotation (stochastic remnant) & roughly flat near $r_*=\sqrt{2GM/|\gamma_1|}$, $v_*^4=2GM|\gamma_1|$, $|\gamma_1|\sim c^2\sqrt\Lambda$ &
dark matter halo needed & preliminary~\cite{oppenheimrusso}; lensing open \\
Contained remnants & primordial black holes of $10^{14}$--$10^{19}$\,kg; Hawking glow at the light edge & no statement & window
open~\cite{gortongreen} \\
Spacetime noise given gravitational decoherence & on the trade-off line~\eqref{eq:saturate} & not required & neither yet measured \\
Decoherence of mesoscopic superpositions (Bridge~B) & $\lambda=GM^2/\hbar R$: 72/s for the $10^{-14}$\,kg microsphere & none & untested at
this scale \\
\nw{\vtag{} Proton decay (a unification before the start)} & \nw{possible; $\tau\sim5\times10^{35}$\,yr for $M_X=10^{16}$\,GeV} & \nw{---} &
\nw{$\tau>2.4\times10^{34}$\,yr~\cite{superk}; Hyper-Kamiokande} \\
\nw{\vtag{} Relic magnetic monopoles} & \nw{none: the strong split lies within the start} & \nw{---} & \nw{flux below $1.4\times10^{-16}$ per
cm$^2$\,s\,sr~\cite{macro}} \\
\bottomrule
\end{longtable}
\end{center}

\section{The ledger}\label{sec:ledger}

Every statement of the account, abridged, with its verdict.

\begin{center}\small
\begin{longtable}{@{}p{0.5\linewidth}p{0.24\linewidth}p{0.18\linewidth}@{}}
\toprule
Statement & Verdict & Where \\
\midrule
\endhead
Free fall is universal for the existence frame & derived; holds & Prop.~\ref{pr:wep}; \chk{Q6} \\
The information scalar does not enter gravity; $\gamma=1$ & derived; holds & Prop.~\ref{pr:gamma}; \chk{Q7} \\
The existence frame is classical after the first record & derived & Lemma~\ref{lem:classical}; \chk{Q3} \\
A classical frame cannot mediate entanglement & derived & Lemma~\ref{lem:locc}; \chk{Q4} \\
DMP, with the base configuration of our universe, hosts the proposed theory & derived & Theorem~\ref{th:host} \\
No graviton as a quantum of the frame & derived; consistent with all data & \chk{Q9} \\
The frame must diffuse and decohere matter & derived (with CQ dynamics) & Cor.~\ref{co:diffuse} \\
Record and recompute is measurement and feedback & derived (structure) & Prop.~\ref{pr:ktm} \\
The balance holds the trade-off at equality & live & \eqref{eq:saturate} \\
Bridge A: the outermost horizon fixes the stochastic remnant & refuted (8 in $v^4$) & \chk{Q19} \\
Bridge B: a record at Planck strength gives the Di\'osi--Penrose rate & live; the BMV masses decohere in 14\,ms & \chk{Q19} \\
Gravity entangles no two masses & live, falsifiable & \chk{Q5} \\
The existence of a boundary $gh/c$ is $4\pi^2kT_U$ & holds as a reading & Prop.~\ref{pr:unruh}; \chk{Q1} \\
Time slows where existence is strong (lapse) & holds & \chk{Q2} \\
Einstein's equation as an equation of state & holds as a reading; constant open & \eqref{eq:jacobson} \\
Entropy as the cost of converting existence and information & holds as a reading (Landauer) & Section~\ref{sec:thermo} \\
Horizon as containment; counts as $M,J,Q$ & holds as a reading & Section~\ref{sec:bh} \\
Christodoulou's mass formula from the containment split & holds as a rederivation & \eqref{eq:christo}; \chk{Q10} \\
Spinning holes give up energy by dragging the frame, outside the horizon & holds & \chk{Q10} \\
Hawking--Penrose--Information radiation & holds as a name; Hawking part untested in nature & \chk{Q11} \\
Information returns from black holes & open & Section~\ref{sec:bh} \\
Holographic minimal cuts and reconstruction from the axioms & holds, on a tree & \chk{A1}--\chk{A5} \\
The area is a number (one fixed-area sector) & derived & \eqref{eq:harlow} \\
The tree is our space & refuted as literal; holds as routing or histories & Section~\ref{sec:holo} \\
2+1 gravity is an exception & refuted; runs as content, not as the frame & Section~\ref{sec:21}; \chk{Q8} \\
Inflation and oscillon decay are one start, the minting & holds as a reading & Hyp.~\ref{h:minting} \\
Quantum geometry only within the start & derived (with Hyp.~\ref{h:minting}) & Section~\ref{sec:cosmo} \\
The early universe a Kugelblitz & holds as a reading, at every epoch & \chk{Q12} \\
Inflation as the possibility space; flatness & holds as a reading & \chk{Q13} \\
Primordial tensor modes, generated within the start & live; consistent with all bounds & \chk{Q14} \\
The balanced state as the de Sitter future & holds as a reading in $\Lambda$CDM & \chk{Q13} \\
Remnants of interactions are trapped in the gravitational field & live & Hyp.~\ref{h:remnant} \\
Stochastic remnant: flat rotation curves, $v^4\propto M$, scale tied to $\Lambda$ & holds in a preliminary calculation; lensing open & \chk{Q18} \\
Contained remnant: primordial black holes from the start & live (window open) & \chk{Q18} \\
Radiated remnant (gravitational waves) as the dark matter & refuted (about $10^{-8}$) & \chk{Q18} \\
The known neutrinos as the remnant & refuted (four bounds) & \chk{Q15} \\
Baryon content fixed early & holds & Section~\ref{sec:cosmo} \\
$\alpha$ does not grow & live & Section~\ref{sec:alpha} \\
\nw{\vtag{} $\alpha$ was never measurably one; it was born near $1/128$ at the electroweak split} & \nw{holds (replaces ``$\alpha$ began at
one'')} & \nw{Section~\ref{sec:alpha}; \chk{Q20}} \\
\nw{\vtag{} Gravity is a separate plane, shaped by the counts of existence and information} & \nw{holds (general relativity; it weighs
all energy)} & \nw{Hyp.~1$'$; \chk{Q6}} \\
\nw{\vtag{} The plane's link to G0 is the causal chain of whatever exists} & \nw{holds as a reading} & \nw{Hyp.~1$'$} \\
\nw{\vtag{} The Higgs vacuum is the existence that separated out at the electroweak split} & \nw{holds as a reading} &
\nw{Section~\ref{sec:higgsplit}; \chk{Q20}} \\
\nw{\vtag{} Before the split, existence through interaction (thermal masses)} & \nw{holds as a reading} & \nw{Section~\ref{sec:higgsplit}} \\
\nw{\vtag{} The symmetric electroweak state became unstable} & \nw{holds (a smooth crossover at 160\,GeV)} & \nw{\chk{Q20}} \\
\nw{\vtag{} The information field split out of the electroweak field} & \nw{holds} & \nw{Section~\ref{sec:higgsplit}; \chk{Q20}} \\
\nw{\vtag{} Causality begins with the strong split; a unification lies before time} & \nw{holds as a reading; the order is required by
the absence of monopoles; unification untested} & \nw{Hyp.~6$'$; \chk{Q20}} \\
\nw{\vtag{} Higgs inflation: the Higgs field drives the start} & \nw{open; in tension with Prop.~\ref{pr:gamma}} &
\nw{Section~\ref{sec:higgsplit}} \\
Spin ladder & holds as a reading & Section~\ref{sec:matter} \\
The Higgs vacuum is G0's grant of rest mass & holds as a reading & Section~\ref{sec:matter} \\
The Higgs field is the gravitational field & refuted & \chk{Q6} \\
$\alpha$ as the VIBE coefficient, self-interaction strength & holds as a reading & \chk{Q16} \\
The forces as tiers; their unity & holds as a reading; unity untested & Section~\ref{sec:matter} \\
\bottomrule
\end{longtable}
\end{center}

\section{Where we expect it to break}\label{sec:break}
\begin{enumerate}[nosep]
\item \emph{The hypotheses.} Everything physical rests on Hypotheses~\ref{h:exist}--\ref{h:remnant}, above all on reading trust as existence
and existence as energy-momentum. Another faithful reading of either changes the propositions; the dissent of Section~\ref{sec:holo} is one,
and it predicts $\mathcal N>0$.
\item \emph{Experiment.} A positive test of gravitationally mediated entanglement, a violation of the LOCC inequality, or decoherence and
diffusion both measured below the trade-off~\eqref{eq:tradeoff} would each show that gravity is not GUID's classical frame. Unless classical
gravity can entangle after all~\cite{azizhowl}, a positive test is decisive.
\item \emph{Constants.} The axioms derive none of $G$, $\hbar$, $\Lambda$, the entropy per area $\eta$, or the strength of the
measurement-and-feedback channel; they enter as content or configuration. The two bridges of Section~\ref{sec:bridges} fix the strength
only by adding a principle, and give a refuted value (A) or the Di\'osi--Penrose rate (B). If mesoscopic superpositions are seen to keep
their coherence well beyond $\hbar R/GM^2$, Bridge~B is refuted too.
\item \emph{The start.} Hypothesis~\ref{h:minting} makes inflation and the oscillon decay one start that the axioms do not constrain, so GUID
predicts nothing about primordial tensor modes beyond standard inflation; its predictions about the geometry begin after the start.
\item \emph{Dark matter.} The stochastic remnant has been computed only for static, spherically symmetric configurations. If it cannot
reproduce lensing, the acoustic peaks and the merging cluster 1E0657$-$558, it is not the dark matter and the contained remnant must carry
it; if evaporation and microlensing close the asteroid-mass window, that route closes too.
\item \emph{Black holes.} GUID does not say whether information returns from an evaporating hole. A derivation of the Page curve that needs
superposed geometries after the start would contradict Lemma~\ref{lem:classical}.
\item \emph{Space.} A literal tree is not our space; the routing reading avoids the objection but has not been shown to give a Lorentzian
3+1 manifold in any limit.
\item \nw{\vtag{} \emph{The plane.} Hypothesis~1$'$ keeps the classical prediction only because the plane is a DMP structure shaped by
the counts. If the plane is read as lying outside DMP's structures, DMP says nothing about whether gravity is quantum.}
\item \nw{\vtag{} \emph{Higgs inflation.} If the Higgs field drives the start, its coupling to curvature must be reconciled with
Proposition~\ref{pr:gamma}; if Axiom~1 forbids it in every form, Higgs inflation and GUID exclude each other.}
\item \nw{\vtag{} \emph{Unification.} A proton decay found would show that a unification existed, as Hypothesis~6$'$ allows; continued
silence at ever longer lifetimes makes a unification less likely but cannot refute one that lies before time.}
\item \emph{The checks.} \chk{Q6} uses the semi-empirical Coulomb term and a coarse composition of the Earth; \chk{Q13} a reheating
temperature of $10^{15}$\,GeV; \chk{Q15} round values for a dwarf galaxy; \chk{Q16} one-loop running. Better inputs move numbers by tens of
per cent, not the verdicts.
\end{enumerate}

\section{Conclusion}

Read as physics, the three axioms of the Dark Matter Protocol give a definite account of quantum gravity, and it is not a quantisation of
geometry. Trust becomes existence, and existence becomes a tetrad that weighs all energy alike, never mixes with the information half, and
is classical once the start, inflation with the decay of its oscillons, is complete, because the axioms make each of its values a record
that two independent witnesses hold in full. Supplied with the base configuration currently assumed for our universe, DMP hosts this
theory: its shared structure carries the classical frame, and its information half the quantum content. Consistency with quantum
mechanics then fixes the rest of the dynamics: the frame must diffuse and decohere matter, and the axioms' own architecture, record then
recompute, is the measurement-and-feedback channel that does so. Horizons are containing boundaries whose temperature is the existence they
assert; black holes give up what the first law allows, and an outside reader sees only their counts; holographic codes are reproduced up to
their classical centre; 2+1 dimensions are no exception; quantum geometry is confined to the start. The theory lets remnants of interactions
be trapped in the gravitational field: its stochastic remnant already gives the shape and scale of galactic rotation curves in a first
calculation, and its contained remnants are primordial black holes from the start. Made parameter-free with its own
principles, it gives either a value the rotation curves refute or the Di\'osi--Penrose decoherence rate, so its physics is that of the
classical-quantum family; what it adds is a reason for that family, drawn from an axiom system for trust. It stands on the classical side
of the one class of experiments that can decide it, gravitationally mediated entanglement and its relatives. If those experiments find entanglement, the
geometry is quantum and the proposed theory is refuted, and DMP would run its successor as content.
\nw{\vtag{} In version 5.1 gravity is the separate plane that existence and information shape, the Higgs vacuum is the existence that
separated out when the information field split from the electroweak field, $\alpha$ was never measurably one, and causality begins with
the split of the strong force from the electroweak force, an order the absence of magnetic monopoles requires.} We claim that the
construction exists
and reproduces. Whether it is right is for the experiments, and for the reader, to break.

\section*{Acknowledgements}
The first author thanks Sir Peter Higgs and the many theorists and experimentalists whose work this account requires and builds on. The
review passes that tested earlier forms of this work, including the dissenting one reported in Section~\ref{sec:holo}, are part of its
record.

\section*{Contributions and use of AI}
Steven Hrovat conceived DMP, the reading of trust as existence, the two-field structure, and every hypothesis tested here, chose which ideas
to keep after each round of tests, and takes responsibility for the axiom system and the implications of its theorems. Sirius (Claude,
Anthropic) made the derivations, related the account to the literature, wrote and ran \texttt{docs/quantum-gravity/checks.py}, and wrote
this text. The holographic-code assembly, the Lean development and the review passes were written by AI systems under the first author's
direction. Every numerical claim is a line of a public script; every reference was checked against a publisher, arXiv or collaboration
record, as listed in \texttt{docs/quantum-gravity/README.md}.

\section*{Licence and availability}
This paper and its source are released under the Apache License, Version 2.0. DMP: \url{https://github.com/StevenHrovat/dmp} (version 0.9.9
at commit \texttt{4400cce}). Gravity: \url{https://github.com/StevenHrovat/gravity} (this paper and its checks in
\texttt{docs/quantum-gravity/}; the holographic-code assembly in \texttt{docs/qg-verdict/}; the Lean development in \texttt{evidence/}).

\appendix
\section{The checks}\label{app:checks}

\texttt{python3 docs/quantum-gravity/checks.py} (numpy, scipy, sympy) prints \nw{123} checks, all of which pass. Q1: Planck units, the existence
energy, Unruh, Hawking, Bekenstein--Hawking, Jacobson's entropy density, the size of the quantum term of~\eqref{eq:eft}. Q2: the lapse. Q3:
no-cloning and broadcasting. Q4: classical against coherent mediation, 300 random couplings each. Q5: the Bose et al.\ geometry. Q6: the
MICROSCOPE bound on the information scalar and the Higgs-share test. Q7: Brans--Dicke and Cassini. Q8: polarisations and Weyl components
in $D$ dimensions. Q9: the single-graviton threshold. Q10: Kerr horizon, ergosphere, Christodoulou, efficiencies. Q11: Hawking radiation of
real holes. Q12: the Hubble mass at its Schwarzschild threshold. Q13: flatness, de Sitter floor, onset of acceleration. Q14: the tensor
spectrum, the BICEP/Keck bound, Starobinsky's $r$. Q15: what every trapped remnant must satisfy, and why the known neutrinos do not. Q16: $\alpha$ as self-interaction strength,
the electron's lengths, gravity against electromagnetism, the running of $\alpha$. Q17: the speed of gravity. Q18: remnants trapped in the field: the stochastic remnant's rotation curve and
acceleration scale, the contained remnant's Hawking temperatures and evaporation, the radiated remnant's energy density. Q19: the two bridges that fix the free constant, Bridge A against the Tully--Fisher
normalisation and Bridge B's windows and Di\'osi--Penrose rates. \nw{\vtag{} Q20: the radiation-era clock and the time of
the electroweak crossover, the W and Z masses from the Higgs vacuum, $\alpha$ at its birth and today, and the proton lifetime for a
unification at the start.}
\texttt{python3 docs/qg-verdict/assembly.py} (numpy) prints the 18 checks of the holographic-code assembly; A1 to A5 in this paper are its
checks [1] to [5], and its check [6] is the special case of \chk{Q4} with CNOT and CZ couplings.

\begin{thebibliography}{999}\small
\bibitem{goroff} M.~H. Goroff, A. Sagnotti, ``The ultraviolet behavior of Einstein gravity'', \emph{Nucl. Phys. B} 266, 709 (1986).
\bibitem{thooftveltman} G. 't~Hooft, M. Veltman, ``One-loop divergencies in the theory of gravitation'', \emph{Ann. Inst. Henri
Poincar\'e A} 20, 69 (1974).
\bibitem{donoghue} J.~F. Donoghue, ``General relativity as an effective field theory: the leading quantum corrections'', \emph{Phys. Rev.
D} 50, 3874 (1994).
\bibitem{bdh} N.~E.~J. Bjerrum-Bohr, J.~F. Donoghue, B.~R. Holstein, ``Quantum gravitational corrections to the nonrelativistic
scattering potential of two masses'', \emph{Phys. Rev. D} 67, 084033 (2003).
\bibitem{hawking75} S.~W. Hawking, ``Particle creation by black holes'', \emph{Commun. Math. Phys.} 43, 199 (1975).
\bibitem{bekenstein73} J.~D. Bekenstein, ``Black holes and entropy'', \emph{Phys. Rev. D} 7, 2333 (1973).
\bibitem{page93} D.~N. Page, ``Information in black hole radiation'', \emph{Phys. Rev. Lett.} 71, 3743 (1993).
\bibitem{penington} G. Penington, ``Entanglement wedge reconstruction and the information paradox'', \emph{JHEP} 09 (2020) 002.
\bibitem{aemm} A. Almheiri, N. Engelhardt, D. Marolf, H. Maxfield, ``The entropy of bulk quantum fields and the entanglement wedge of an
evaporating black hole'', \emph{JHEP} 12 (2019) 063.
\bibitem{eppley} K. Eppley, E. Hannah, ``The necessity of quantizing the gravitational field'', \emph{Found. Phys.} 7, 51 (1977).
\bibitem{pagegeilker} D.~N. Page, C.~D. Geilker, ``Indirect evidence for quantum gravity'', \emph{Phys. Rev. Lett.} 47, 979 (1981).
\bibitem{oppenheim} J. Oppenheim, ``A postquantum theory of classical gravity?'', \emph{Phys. Rev. X} 13, 041040 (2023).
\bibitem{osswd} J. Oppenheim, C. Sparaciari, B. \v{S}oda, Z. Weller-Davies, ``Gravitationally induced decoherence vs space-time diffusion:
testing the quantum nature of gravity'', \emph{Nat. Commun.} 14, 7910 (2023), arXiv:2203.01982.
\bibitem{bose} S. Bose, A. Mazumdar, G.~W. Morley, H. Ulbricht, M. Toro\v{s}, M. Paternostro, A.~A. Geraci, P.~F. Barker, M.~S. Kim, G.
Milburn, ``Spin entanglement witness for quantum gravity'', \emph{Phys. Rev. Lett.} 119, 240401 (2017).
\bibitem{mv} C. Marletto, V. Vedral, ``Gravitationally induced entanglement between two massive particles is sufficient evidence of
quantum effects in gravity'', \emph{Phys. Rev. Lett.} 119, 240402 (2017).
\bibitem{ggs} T.~D. Galley, F. Giacomini, J.~H. Selby, ``A no-go theorem on the nature of the gravitational field beyond quantum
theory'', \emph{Quantum} 6, 779 (2022).
\bibitem{dmp} S. Hrovat, \emph{The Dark Matter Protocol}, version 0.9.9, \url{https://github.com/StevenHrovat/dmp}, commit
\texttt{4400cce} (2026).
\bibitem{gm} J.~A. Goguen, J. Meseguer, ``Security policies and security models'', in \emph{IEEE Symposium on Security and Privacy}, 11
(1982).
\bibitem{denning} D.~E. Denning, ``A lattice model of secure information flow'', \emph{Commun. ACM} 19, 236 (1976).
\bibitem{mr} D. Malkhi, M. Reiter, ``Byzantine quorum systems'', \emph{Distrib. Comput.} 11, 203 (1998).
\bibitem{ohearn} P. O'Hearn, J. Reynolds, H. Yang, ``Local reasoning about programs that alter data structures'', in \emph{CSL 2001},
LNCS 2142, 1 (2001).
\bibitem{lean4} L. de~Moura, S. Ullrich, ``The Lean 4 theorem prover and programming language'', in \emph{CADE 28}, LNCS 12699, 625
(2021).
\bibitem{aldrovandi} R. Aldrovandi, J.~G. Pereira, \emph{Teleparallel Gravity: An Introduction} (Springer, Dordrecht, 2013).
\bibitem{aether} T. Jacobson, D. Mattingly, ``Gravity with a dynamical preferred frame'', \emph{Phys. Rev. D} 64, 024028 (2001).
\bibitem{gw170817} B.~P. Abbott et al. (LIGO, Virgo, Fermi GBM, INTEGRAL), ``Gravitational waves and gamma-rays from a binary neutron
star merger: GW170817 and GRB 170817A'', \emph{Astrophys. J. Lett.} 848, L13 (2017).
\bibitem{bekenstein82} J.~D. Bekenstein, ``Fine-structure constant: Is it really a constant?'', \emph{Phys. Rev. D} 25, 1527 (1982).
\bibitem{bsbm} H.~B. Sandvik, J.~D. Barrow, J. Magueijo, ``A simple cosmology with a varying fine structure constant'', \emph{Phys. Rev.
Lett.} 88, 031302 (2002).
\bibitem{microscope} P. Touboul et al. (MICROSCOPE Collaboration), ``MICROSCOPE mission: final results of the test of the equivalence
principle'', \emph{Phys. Rev. Lett.} 129, 121102 (2022).
\bibitem{yang} Y.-B. Yang, J. Liang, Y.-J. Bi, Y. Chen, T. Draper, K.-F. Liu, Z. Liu, ``Proton mass decomposition from the QCD energy
momentum tensor'', \emph{Phys. Rev. Lett.} 121, 212001 (2018).
\bibitem{dz} G. Dvali, M. Zaldarriaga, ``Changing $\alpha$ with time: implications for fifth-force-type experiments and quintessence'',
\emph{Phys. Rev. Lett.} 88, 091303 (2002).
\bibitem{will} C.~M. Will, ``The confrontation between general relativity and experiment'', \emph{Living Rev. Relativ.} 17, 4 (2014).
\bibitem{bd} C. Brans, R.~H. Dicke, ``Mach's principle and a relativistic theory of gravitation'', \emph{Phys. Rev.} 124, 925 (1961).
\bibitem{cassini} B. Bertotti, L. Iess, P. Tortora, ``A test of general relativity using radio links with the Cassini spacecraft'',
\emph{Nature} 425, 374 (2003).
\bibitem{wz} W.~K. Wootters, W.~H. Zurek, ``A single quantum cannot be cloned'', \emph{Nature} 299, 802 (1982).
\bibitem{dieks} D. Dieks, ``Communication by EPR devices'', \emph{Phys. Lett. A} 92, 271 (1982).
\bibitem{barnum} H. Barnum, C.~M. Caves, C.~A. Fuchs, R. Jozsa, B. Schumacher, ``Noncommuting mixed states cannot be broadcast'',
\emph{Phys. Rev. Lett.} 76, 2818 (1996).
\bibitem{zurek} W.~H. Zurek, ``Quantum Darwinism'', \emph{Nat. Phys.} 5, 181 (2009).
\bibitem{bph} F.~G.~S.~L. Brand\~ao, M. Piani, P. Horodecki, ``Generic emergence of classical features in quantum Darwinism'',
\emph{Nat. Commun.} 6, 7908 (2015).
\bibitem{vw} G. Vidal, R.~F. Werner, ``Computable measure of entanglement'', \emph{Phys. Rev. A} 65, 032314 (2002).
\bibitem{deutsch} D. Deutsch, ``Quantum theory, the Church--Turing principle and the universal quantum computer'', \emph{Proc. R. Soc.
Lond. A} 400, 97 (1985).
\bibitem{weinberg64} S. Weinberg, ``Photons and gravitons in S-matrix theory: derivation of charge conservation and equality of
gravitational and inertial mass'', \emph{Phys. Rev.} 135, B1049 (1964).
\bibitem{deser70} S. Deser, ``Self-interaction and gauge invariance'', \emph{Gen. Relativ. Gravit.} 1, 9 (1970).
\bibitem{dyson} F. Dyson, ``Is a graviton detectable?'', \emph{Int. J. Mod. Phys. A} 28, 1330041 (2013).
\bibitem{rb} T. Rothman, S. Boughn, ``Can gravitons be detected?'', \emph{Found. Phys.} 36, 1801 (2006).
\bibitem{mead} C.~A. Mead, ``Possible connection between gravitation and fundamental length'', \emph{Phys. Rev.} 135, B849 (1964).
\bibitem{tobar} G. Tobar, S.~K. Manikandan, T. Beitel, I. Pikovski, ``Detecting single gravitons with quantum sensing'', \emph{Nat.
Commun.} 15, 7229 (2024).
\bibitem{cdr} D. Carney, V. Domcke, N.~L. Rodd, ``Graviton detection and the quantization of gravity'', \emph{Phys. Rev. D} 109, 044009
(2024).
\bibitem{moller} C. M{\o}ller, ``The energy-momentum complex in general relativity and related problems'', in \emph{Les Th\'eories
Relativistes de la Gravitation}, ed. A. Lichnerowicz, M.-A. Tonnelat (CNRS, Paris, 1962).
\bibitem{rosenfeld63} L. Rosenfeld, ``On quantization of fields'', \emph{Nucl. Phys.} 40, 353 (1963).
\bibitem{ktm} D. Kafri, J.~M. Taylor, G.~J. Milburn, ``A classical channel model for gravitational decoherence'', \emph{New J. Phys.} 16,
065020 (2014).
\bibitem{tilloydiosi} A. Tilloy, L. Di\'osi, ``Sourcing semiclassical gravity from spontaneously localized quantum matter'', \emph{Phys.
Rev. D} 93, 024026 (2016).
\bibitem{diosi87} L. Di\'osi, ``A universal master equation for the gravitational violation of quantum mechanics'', \emph{Phys. Lett. A}
120, 377 (1987).
\bibitem{penrose96} R. Penrose, ``On gravity's role in quantum state reduction'', \emph{Gen. Relativ. Gravit.} 28, 581 (1996).
\bibitem{donadi} S. Donadi, K. Piscicchia, C. Curceanu, L. Di\'osi, M. Laubenstein, A. Bassi, ``Underground test of gravity-related wave
function collapse'', \emph{Nat. Phys.} 17, 74 (2021).
\bibitem{azizhowl} J. Aziz, R. Howl, ``Classical theories of gravity produce entanglement'', \emph{Nature} 646, 813 (2025).
\bibitem{movw} C. Marletto, J. Oppenheim, V. Vedral, E. Wilson, ``Classical gravity cannot mediate entanglement'', arXiv:2511.07348
(2025).
\bibitem{diosi25} L. Di\'osi, ``No, classical gravity does not entangle quantized matter fields'', arXiv:2511.00852 (2025).
\bibitem{shl} M.~D. Schneider, N. Huggett, N. Linnemann, ``A demonstration that classical gravity does not produce entanglement'',
arXiv:2511.19242 (2025).
\bibitem{gib} L. Gundhi, A. Infantino, A. Bassi, ``Can classical theories of gravity produce entanglement?'', arXiv:2604.19696 (2026).
\bibitem{christodoulou23} M. Christodoulou, A. Di~Biagio, M. Aspelmeyer, \v{C}. Brukner, C. Rovelli, R. Howl, ``Locally mediated
entanglement in linearized quantum gravity'', \emph{Phys. Rev. Lett.} 130, 100202 (2023).
\bibitem{lpp} L. Lami, J.~S. Pedernales, M.~B. Plenio, ``Testing the quantumness of gravity without entanglement'', \emph{Phys. Rev. X} 14,
021022 (2024).
\bibitem{westphal} T. Westphal, H. Hepach, J. Pfaff, M. Aspelmeyer, ``Measurement of gravitational coupling between millimetre-sized
masses'', \emph{Nature} 591, 225 (2021).
\bibitem{cow} R. Colella, A.~W. Overhauser, S.~A. Werner, ``Observation of gravitationally induced quantum interference'', \emph{Phys. Rev.
Lett.} 34, 1472 (1975).
\bibitem{nesvizhevsky} V.~V. Nesvizhevsky et al., ``Quantum states of neutrons in the Earth's gravitational field'', \emph{Nature} 415,
297 (2002).
\bibitem{cst} D. Carney, P.~C.~E. Stamp, J.~M. Taylor, ``Tabletop experiments for quantum gravity: a user's manual'', \emph{Class. Quantum
Grav.} 36, 034001 (2019).
\bibitem{unruh} W.~G. Unruh, ``Notes on black-hole evaporation'', \emph{Phys. Rev. D} 14, 870 (1976).
\bibitem{pr60} R.~V. Pound, G.~A. Rebka, Jr., ``Apparent weight of photons'', \emph{Phys. Rev. Lett.} 4, 337 (1960).
\bibitem{chou} C.~W. Chou, D.~B. Hume, T. Rosenband, D.~J. Wineland, ``Optical clocks and relativity'', \emph{Science} 329, 1630 (2010).
\bibitem{jacobson95} T. Jacobson, ``Thermodynamics of spacetime: the Einstein equation of state'', \emph{Phys. Rev. Lett.} 75, 1260 (1995).
\bibitem{jacobson16} T. Jacobson, ``Entanglement equilibrium and the Einstein equation'', \emph{Phys. Rev. Lett.} 116, 201101 (2016).
\bibitem{landauer} R. Landauer, ``Irreversibility and heat generation in the computing process'', \emph{IBM J. Res. Dev.} 5, 183 (1961).
\bibitem{israel} W. Israel, ``Event horizons in static vacuum space-times'', \emph{Phys. Rev.} 164, 1776 (1967).
\bibitem{kerr} R.~P. Kerr, ``Gravitational field of a spinning mass as an example of algebraically special metrics'', \emph{Phys. Rev.
Lett.} 11, 237 (1963).
\bibitem{christodoulou} D. Christodoulou, ``Reversible and irreversible transformations in black-hole physics'', \emph{Phys. Rev. Lett.}
25, 1596 (1970).
\bibitem{bch} J.~M. Bardeen, B. Carter, S.~W. Hawking, ``The four laws of black hole mechanics'', \emph{Commun. Math. Phys.} 31, 161
(1973).
\bibitem{penrose69} R. Penrose, ``Gravitational collapse: the role of general relativity'', \emph{Riv. Nuovo Cimento} 1, 252 (1969).
\bibitem{bz} R.~D. Blandford, R.~L. Znajek, ``Electromagnetic extraction of energy from Kerr black holes'', \emph{Mon. Not. R. Astron.
Soc.} 179, 433 (1977).
\bibitem{bpt} J.~M. Bardeen, W.~H. Press, S.~A. Teukolsky, ``Rotating black holes: locally nonrotating frames, energy extraction, and
scalar synchrotron radiation'', \emph{Astrophys. J.} 178, 347 (1972).
\bibitem{happy} F. Pastawski, B. Yoshida, D. Harlow, J. Preskill, ``Holographic quantum error-correcting codes: toy models for the
bulk/boundary correspondence'', \emph{JHEP} 06 (2015) 149.
\bibitem{laflamme} R. Laflamme, C. Miquel, J.~P. Paz, W.~H. Zurek, ``Perfect quantum error correcting code'', \emph{Phys. Rev. Lett.}
77, 198 (1996).
\bibitem{rt} S. Ryu, T. Takayanagi, ``Holographic derivation of entanglement entropy from AdS/CFT'', \emph{Phys. Rev. Lett.} 96, 181602
(2006).
\bibitem{cgl} R. Cleve, D. Gottesman, H.-K. Lo, ``How to share a quantum secret'', \emph{Phys. Rev. Lett.} 83, 648 (1999).
\bibitem{adh} A. Almheiri, X. Dong, D. Harlow, ``Bulk locality and quantum error correction in AdS/CFT'', \emph{JHEP} 04 (2015) 163.
\bibitem{harlow} D. Harlow, ``The Ryu--Takayanagi formula from quantum error correction'', \emph{Commun. Math. Phys.} 354, 865 (2017).
\bibitem{dhm} X. Dong, D. Harlow, D. Marolf, ``Flat entanglement spectra in fixed-area states of quantum gravity'', \emph{JHEP} 10
(2019) 240.
\bibitem{hmss} M. Heydeman, M. Marcolli, I. Saberi, B. Stoica, ``Tensor networks, $p$-adic fields, and algebraic curves: arithmetic and
the AdS$_3$/CFT$_2$ correspondence'', \emph{Adv. Theor. Math. Phys.} 22, 93 (2018).
\bibitem{gubser} S.~S. Gubser, J. Knaute, S. Parikh, A. Samberg, P. Witaszczyk, ``$p$-adic AdS/CFT'', \emph{Commun. Math. Phys.} 352,
1019 (2017).
\bibitem{rr} Y. Rabinovich, R. Raz, ``Lower bounds on the distortion of embedding finite metric spaces in graphs'', \emph{Discrete
Comput. Geom.} 19, 79 (1998).
\bibitem{gellmannhartle} M. Gell-Mann, J.~B. Hartle, ``Quantum mechanics in the light of quantum cosmology'', in \emph{Complexity, Entropy
and the Physics of Information}, ed. W.~H. Zurek (Addison-Wesley, Redwood City, 1990), 425.
\bibitem{djt} S. Deser, R. Jackiw, G. 't~Hooft, ``Three-dimensional Einstein gravity: dynamics of flat space'', \emph{Ann. Phys.} 152, 220
(1984).
\bibitem{witten88} E. Witten, ``2+1 dimensional gravity as an exactly soluble system'', \emph{Nucl. Phys. B} 311, 46 (1988).
\bibitem{bh} J.~D. Brown, M. Henneaux, ``Central charges in the canonical realization of asymptotic symmetries: an example from
three-dimensional gravity'', \emph{Commun. Math. Phys.} 104, 207 (1986).
\bibitem{mw} A. Maloney, E. Witten, ``Quantum gravity partition functions in three dimensions'', \emph{JHEP} 02 (2010) 029.
\bibitem{nr} J.~E. Nelson, T. Regge, ``Homotopy groups and (2+1)-dimensional quantum gravity'', \emph{Nucl. Phys. B} 328, 190 (1989).
\bibitem{thooft88} G. 't~Hooft, ``Non-perturbative 2 particle scattering amplitudes in 2+1 dimensional quantum gravity'',
\emph{Commun. Math. Phys.} 117, 685 (1988).
\bibitem{gw150914} B.~P. Abbott et al. (LIGO Scientific and Virgo Collaborations), ``Observation of gravitational waves from a binary black
hole merger'', \emph{Phys. Rev. Lett.} 116, 061102 (2016).
\bibitem{halliwell} J.~J. Halliwell, ``Decoherence in quantum cosmology'', \emph{Phys. Rev. D} 39, 2912 (1989).
\bibitem{kiefer87} C. Kiefer, ``Continuous measurement of mini-superspace variables by higher multipoles'', \emph{Class. Quantum Grav.} 4,
1369 (1987).
\bibitem{hartlehawking} J.~B. Hartle, S.~W. Hawking, ``Wave function of the Universe'', \emph{Phys. Rev. D} 28, 2960 (1983).
\bibitem{pathria} R.~K. Pathria, ``The universe as a black hole'', \emph{Nature} 240, 298 (1972).
\bibitem{guth} A.~H. Guth, ``Inflationary universe: a possible solution to the horizon and flatness problems'', \emph{Phys. Rev. D} 23, 347
(1981).
\bibitem{planck18} Planck Collaboration, ``Planck 2018 results. VI. Cosmological parameters'', \emph{Astron. Astrophys.} 641, A6 (2020).
\bibitem{starobinsky79} A.~A. Starobinsky, ``Spectrum of relict gravitational radiation and the early state of the universe'', \emph{JETP
Lett.} 30, 682 (1979).
\bibitem{krausswilczek} L.~M. Krauss, F. Wilczek, ``Using cosmology to establish the quantization of gravity'', \emph{Phys. Rev. D} 89,
047501 (2014).
\bibitem{bk18} P.~A.~R. Ade et al. (BICEP/Keck), ``Improved constraints on primordial gravitational waves using Planck, WMAP, and BICEP/Keck
observations through the 2018 observing season'', \emph{Phys. Rev. Lett.} 127, 151301 (2021).
\bibitem{tristram} M. Tristram et al., ``Improved limits on the tensor-to-scalar ratio using BICEP and Planck data'', \emph{Phys. Rev. D}
105, 083524 (2022).
\bibitem{starobinsky80} A.~A. Starobinsky, ``A new type of isotropic cosmological models without singularity'', \emph{Phys. Lett. B} 91,
99 (1980).
\bibitem{litebird} E. Allys et al. (LiteBIRD Collaboration), ``Probing cosmic inflation with the LiteBIRD cosmic microwave background
polarization survey'', \emph{Prog. Theor. Exp. Phys.} 2023, 042F01 (2023).
\bibitem{polarski} D. Polarski, A.~A. Starobinsky, ``Semiclassicality and decoherence of cosmological perturbations'', \emph{Class. Quantum
Grav.} 13, 377 (1996).
\bibitem{desi2} DESI Collaboration, ``DESI DR2 results II: measurements of baryon acoustic oscillations and cosmological constraints'',
\emph{Phys. Rev. D} 112, 083515 (2025).
\bibitem{desi} DESI Collaboration, ``DESI 2024 VI: cosmological constraints from the measurements of baryon acoustic oscillations'',
\emph{JCAP} 02 (2025) 021.
\bibitem{tremainegunn} S. Tremaine, J.~E. Gunn, ``Dynamical role of light neutral leptons in cosmology'', \emph{Phys. Rev. Lett.} 42, 407
(1979).
\bibitem{piffl} T. Piffl et al., ``The RAVE survey: the Galactic escape speed and the mass of the Milky Way'', \emph{Astron. Astrophys.}
562, A91 (2014).
\bibitem{lange} R. Lange et al., ``Improved limits for violations of local position invariance from atomic clock comparisons'',
\emph{Phys. Rev. Lett.} 126, 011102 (2021).
\bibitem{espresso} M.~T. Murphy et al., ``Fundamental physics with ESPRESSO: precise limit on variations in the fine-structure constant
towards the bright quasar HE 0515$-$4414'', \emph{Astron. Astrophys.} 658, A123 (2022).
\bibitem{planckxxiv} Planck Collaboration, ``Planck intermediate results. XXIV. Constraints on variations in fundamental constants'',
\emph{Astron. Astrophys.} 580, A22 (2015).
\bibitem{higgs64} P.~W. Higgs, ``Broken symmetries and the masses of gauge bosons'', \emph{Phys. Rev. Lett.} 13, 508 (1964).
\bibitem{be64} F. Englert, R. Brout, ``Broken symmetry and the mass of gauge vector mesons'', \emph{Phys. Rev. Lett.} 13, 321 (1964).
\bibitem{schwinger} J. Schwinger, ``On quantum-electrodynamics and the magnetic moment of the electron'', \emph{Phys. Rev.} 73, 416
(1948).
\bibitem{fan} X. Fan, T.~G. Myers, B.~A.~D. Sukra, G. Gabrielse, ``Measurement of the electron magnetic moment'', \emph{Phys. Rev. Lett.}
130, 071801 (2023).
\bibitem{glashow} S.~L. Glashow, ``Partial-symmetries of weak interactions'', \emph{Nucl. Phys.} 22, 579 (1961).
\bibitem{weinberg67} S. Weinberg, ``A model of leptons'', \emph{Phys. Rev. Lett.} 19, 1264 (1967).
\bibitem{salam} A. Salam, ``Weak and electromagnetic interactions'', in \emph{Elementary Particle Theory}, ed. N. Svartholm (Almqvist and
Wiksell, Stockholm, 1968), 367.
\bibitem{amaldi} U. Amaldi, W. de~Boer, H. F{\"u}rstenau, ``Comparison of grand unified theories with electroweak and strong coupling
constants measured at LEP'', \emph{Phys. Lett. B} 260, 447 (1991).
\bibitem{maldacena} J.~M. Maldacena, ``The large N limit of superconformal field theories and supergravity'', \emph{Adv. Theor. Math. Phys.}
2, 231 (1998).
\bibitem{ashtekar} A. Ashtekar, ``New variables for classical and quantum gravity'', \emph{Phys. Rev. Lett.} 57, 2244 (1986).
\bibitem{rovellismolin} C. Rovelli, L. Smolin, ``Discreteness of area and volume in quantum gravity'', \emph{Nucl. Phys. B} 442, 593
(1995).
\bibitem{weinberg79} S. Weinberg, ``Ultraviolet divergences in quantum theories of gravitation'', in \emph{General Relativity: An Einstein
Centenary Survey}, ed. S.~W. Hawking, W. Israel (Cambridge University Press, 1979), 790.
\bibitem{reuter} M. Reuter, ``Nonperturbative evolution equation for quantum gravity'', \emph{Phys. Rev. D} 57, 971 (1998).
\bibitem{cdt} J. Ambj{\o}rn, A. G{\"o}rlich, J. Jurkiewicz, R. Loll, ``Planckian birth of the quantum de Sitter universe'',
arXiv:0712.2485 (2007).
\bibitem{bombelli} L. Bombelli, J. Lee, D. Meyer, R.~D. Sorkin, ``Space-time as a causal set'', \emph{Phys. Rev. Lett.} 59, 521 (1987).
\bibitem{vanraamsdonk} M. Van~Raamsdonk, ``Building up spacetime with quantum entanglement'', \emph{Gen. Relativ. Gravit.} 42, 2323 (2010).
\bibitem{ermepr} J. Maldacena, L. Susskind, ``Cool horizons for entangled black holes'', \emph{Fortschr. Phys.} 61, 781 (2013).
\bibitem{oppenheimrusso} J. Oppenheim, A. Russo, ``Anomalous contribution to galactic rotation curves due to stochastic spacetime'',
arXiv:2402.19459 (2024).
\bibitem{milgrom} M. Milgrom, ``A modification of the Newtonian dynamics as a possible alternative to the hidden mass hypothesis'',
\emph{Astrophys. J.} 270, 365 (1983).
\bibitem{hehlmashhoon} F.~W. Hehl, B. Mashhoon, ``Nonlocal gravity simulates dark matter'', arXiv:0812.1059 (2008).
\bibitem{amin} M.~A. Amin, R. Easther, H. Finkel, R. Flauger, M.~P. Hertzberg, ``Oscillons after inflation'', \emph{Phys. Rev. Lett.} 108,
241302 (2012).
\bibitem{kls} L. Kofman, A. Linde, A.~A. Starobinsky, ``Reheating after inflation'', \emph{Phys. Rev. Lett.} 73, 3195 (1994).
\bibitem{cotner} E. Cotner, A. Kusenko, V. Takhistov, ``Primordial black holes from inflaton fragmentation into oscillons'', \emph{Phys.
Rev. D} 98, 083513 (2018).
\bibitem{gortongreen} M. Gorton, A.~M. Green, ``How open is the asteroid-mass primordial black hole window?'', \emph{SciPost Phys.} 17, 032
(2024).
\bibitem{nanograv} G. Agazie et al. (NANOGrav Collaboration), ``The NANOGrav 15 yr data set: evidence for a gravitational-wave
background'', \emph{Astrophys. J. Lett.} 951, L8 (2023).
\bibitem{figurato} L. Figurato, M. Dirindin, J.~L. Gaona-Reyes, M. Carlesso, A. Bassi, S. Donadi, ``On the effectiveness of the collapse
in the Di\'osi--Penrose model'', \emph{New J. Phys.} 26, 113004 (2024).
\bibitem{preskill} J.~P. Preskill, ``Cosmological production of superheavy magnetic monopoles'', \emph{Phys. Rev. Lett.} 43, 1365 (1979).
\bibitem{macro} M. Ambrosio et al. (MACRO Collaboration), ``Final results of magnetic monopole searches with the MACRO experiment'',
\emph{Eur. Phys. J. C} 25, 511 (2002).
\bibitem{superk} A. Takenaka et al. (Super-Kamiokande Collaboration), ``Search for proton decay via $p\to e^+\pi^0$ and $p\to\mu^+\pi^0$
with an enlarged fiducial volume in Super-Kamiokande I-IV'', \emph{Phys. Rev. D} 102, 112011 (2020).
\bibitem{hyperk} K. Abe et al. (Hyper-Kamiokande Proto-Collaboration), ``Hyper-Kamiokande design report'', arXiv:1805.04163 (2018).
\bibitem{georgiglashow} H. Georgi, S.~L. Glashow, ``Unity of all elementary-particle forces'', \emph{Phys. Rev. Lett.} 32, 438 (1974).
\bibitem{donofrio} M. D'Onofrio, K. Rummukainen, ``Standard model cross-over on the lattice'', \emph{Phys. Rev. D} 93, 025003 (2016).
\bibitem{bezrukov} F.~L. Bezrukov, M. Shaposhnikov, ``The Standard Model Higgs boson as the inflaton'', \emph{Phys. Lett. B} 659, 703
(2008).
\bibitem{clowe} D. Clowe, M. Brada\v{c}, A.~H. Gonzalez, M. Markevitch, S.~W. Randall, C. Jones, D. Zaritsky, ``A direct empirical proof of
the existence of dark matter'', \emph{Astrophys. J. Lett.} 648, L109 (2006).
\end{thebibliography}

\end{document}
